% ============================================================
%  Advanced Engineering Informatics (Elsevier) submission
%  Built on the 'cas-dc' double-column CAS Bundle class.
% ============================================================
\documentclass[a4paper,fleqn]{cas-dc}

\usepackage[numbers,sort&compress]{natbib}
\usepackage{graphicx}
\usepackage{amsmath,amssymb,amsfonts,mathtools}
\usepackage{bm}
\usepackage{booktabs}
\usepackage{multirow}
\usepackage{tabularx}
\usepackage{array}
\usepackage{siunitx}
\sisetup{detect-all=true}
\usepackage{enumitem}
\usepackage{adjustbox}
\usepackage{doi}
\usepackage{algorithm}
\usepackage{algpseudocode}

% ---- Hyperlinks ---------------------------------------------
\hypersetup{colorlinks=true,
            linkcolor=black,
            citecolor=black,
            urlcolor=blue}

% ---- Caption spacing ----------------------------------------
\setlength{\abovecaptionskip}{3pt}
\setlength{\belowcaptionskip}{4pt}
\raggedbottom

% ---- Float Placement Controls -------------------------------
\renewcommand{\topfraction}{0.95}
\renewcommand{\bottomfraction}{0.50}
\renewcommand{\textfraction}{0.05}
\renewcommand{\floatpagefraction}{0.85}
\renewcommand{\dbltopfraction}{0.95}
\renewcommand{\dblfloatpagefraction}{0.85}
\setcounter{topnumber}{4}
\setcounter{bottomnumber}{2}
\setcounter{totalnumber}{8}
\setcounter{dbltopnumber}{4}

% ---- Mathematical Notation Macros ---------------------------
\newcommand{\Dcrit}{D_{\mathrm{crit}}}
\newcommand{\RUL}{\mathrm{RUL}}
\newcommand{\MAE}{\mathrm{MAE}}
\newcommand{\RMSE}{\mathrm{RMSE}}
\newcommand{\bPhi}{\bm{\Phi}}
\newcommand{\R}{\mathbb{R}}

\begin{document}
\let\WriteBookmarks\relax
\def\floatpagepagefraction{1}
\def\textpagefraction{.001}

% Short title and short author running head
\shortauthors{H.S. Thota et al.}
\shorttitle{Knowledge-Informed State-Space Degradation Modelling for Ball-Screw RUL}

\title[mode=title]{Knowledge-informed state-space degradation modelling with conformal uncertainty quantification for ball-screw remaining useful life prediction}

\author[1]{Harikrishna S. Thota}[orcid=0000-0002-8646-0428]
\cormark[1]
\ead{harithota@cmti.res.in}
\credit{Conceptualization, Methodology, Software, Formal analysis, Validation, Visualization, Writing -- original draft, Writing -- review \& editing}

\author[2]{P.S. Saikrishna}[orcid=0000-0002-3642-1279]
\credit{Investigation, Validation, Data curation, Supervision, Writing -- review \& editing}

\author[1]{Nagahanumaiah}
\credit{Supervision, Project administration, Resources, Funding acquisition}

\affiliation[1]{organization={Central Manufacturing Technology Institute},
               addressline={Tumkur Road},
               city={Bengaluru},
               postcode={560022},
               state={Karnataka},
               country={India}}

\affiliation[2]{organization={Department of Electrical Engineering, Indian Institute of Technology Tirupati},
               addressline={Yerpedu},
               city={Tirupati},
               postcode={517619},
               state={Andhra Pradesh},
               country={India}}

\cortext[1]{Corresponding author}

% ========================  ABSTRACT  ========================
\begin{abstract}
Remaining useful life (RUL) estimation for ball-screw feed drives is hindered by interacting wear mechanisms, non-stationary duty cycles and unit-to-unit heterogeneity. This paper presents a Knowledge-Informed State-Space Degradation Model (KI-SSDM) that encodes engineering knowledge in four structural elements: (1) a contact-fatigue backbone derived from ISO 3408-5 and Lundberg--Palmgren mechanics; (2) monotone damage accumulation ($\mathrm{d}D/\mathrm{d}t \ge 0$); (3) an Ornstein--Uhlenbeck nuisance process, integrated out analytically, that separates reversible thermal--operational fluctuations from wear; and (4) log-normal unit frailties. Structural discrepancies are learned by a neural or radial-basis correction whose log-rate is bounded by a smooth $\tanh$ map ($e^{-b} \le \alpha \le e^{b}$), and a validation-selected linear pool falls back to pure physics when the correction does not help. Origin-stratified split-conformal prediction provides calibrated intervals. On a 30-fleet synthetic benchmark (2,220 simulated units, 480 test units, six scenarios), KI-SSDM reached fleet-mean MAEs of 16.40~h (planned duty) and 18.29~h (probabilistic duty). After Holm correction it was significantly more accurate than an inverse PINN (planned duty, $p_{\mathrm{Holm}} = 0.0079$) and a deep LSTM (both duties, $p_{\mathrm{Holm}} < 10^{-9}$), and statistically indistinguishable from physics-only, particle-filter and basis-hybrid baselines. Its gain over uncorrected physics was concentrated under load-exponent misspecification (+5.99~h, 29.8\%, planned duty). Calibrated intervals covered 93.3\% and 92.5\% of test cases at mean widths of 82.5~h and 96.4~h. Scheduling maintenance at the conformal lower bound reduced expected cost by 34.7\% (planned) and 31.6\% (probabilistic) relative to periodic $L_{10}$ replacement; the comparators, scheduled the same way, saved 30--35\%.
\end{abstract}

\begin{keywords}
knowledge-informed machine learning \sep engineering informatics \sep remaining useful life \sep state-space models \sep conformal prediction \sep maintenance decision support \sep ball-screw feed drive
\end{keywords}

\maketitle

% ========================  HIGHLIGHTS  ========================
\begin{highlights}
\item Knowledge-informed state-space model for ball-screw remaining useful life.
\item Reversible nuisance state is integrated out exactly; damage stays monotone.
\item Bounded neural correction is used only when validation shows it helps.
\item Split-conformal intervals give 93\% coverage across 30 synthetic fleets.
\item Gain over physics is concentrated under load-exponent misspecification.
\end{highlights}

% ========================  1. INTRODUCTION  ========================
\section{Introduction}
\label{sec:intro}

Precision ball-screw feed drives (BSFDs) in CNC machine tools and robotics operate under multi-axial stress, variable-speed motion, and demanding duty cycles~\cite{Lei_2018, H_Li_2024, Wu_2024, Aivaliotis_2021}. Degradation manifests as rolling-contact fatigue, raceway micro-pitting, lubrication breakdown, and preload loss~\cite{Hanbo_2024, Martin_2024, Heydarnia_2025}. In high-value manufacturing, accurate remaining useful life (RUL) estimation is essential to prevent tool collisions, eliminate premature component overhauls, and enable predictive maintenance~\cite{Lei_2018, H_Li_2024, Zhang_2025, Tai_2023, Yuan_2024, Jia_2019}.

In engineering informatics, condition-based prognostics requires representing degradation through structured, computable engineering knowledge rather than uninterpretable empirical pattern matching~\cite{Liao_2023, D_He_2025, vonRueden_2023}. In mechatronic systems, physical knowledge spans empirical contact fatigue laws (ISO~3408-5)~\cite{ISO_3408_5}, dynamic tribological contact physics (LuGre friction), thermodynamic dissipation, and unit manufacturing variations. Translating this domain knowledge into robust online prognostics is hindered by four core challenges:
\begin{enumerate}[label=(\roman*),leftmargin=2.2em]
    \item \textbf{Decoupling Reversible Dynamics from Irreversible Wear}: Controller-resident telemetry (torque, current, power) reflects high-frequency dynamics and reversible thermal variations ($t_{1/2} \approx 13.9$~h) alongside slow, irreversible fatigue. Disentangling reversible operational nuisance from monotonic degradation without external vibration sensors is challenging~\cite{Pandhare_2021}, and transient operational shifts must not be conflated with unphysical material healing.
    \item \textbf{Physical Inadmissibility in Deep Learning}: Purely data-driven sequence architectures (e.g., LSTMs, transformers) fit training benchmarks well~\cite{Wang_2022, Wei_2024} but lack explicit physical constraints. Nothing in their structure prevents physically inadmissible outputs under unseen duty cycles or data scarcity, such as health indicators that decrease (implied healing) or extrapolations that diverge.
    \item \textbf{Ad-hoc Hybridization and Uncalibrated Uncertainty}: Physics-Informed Neural Networks (PINNs) typically append an empirical $L_2$ ODE residual to a neural loss~\cite{Yan_2026, Gao_2024, Xu_2023, P_Zhang_2026}. This heuristic lacks statistical likelihood semantics, suffers from optimisation gradient imbalances, and yields uncalibrated predictive intervals under distribution shifts~\cite{Xu_2023, Lai_2024, Zhang_CP_2025}.
    \item \textbf{Disconnect from Downstream Decision-Making}: Most prognostic studies terminate at scalar RUL metrics or theoretical error distributions without informing concrete engineering actions. Industrial asset managers require risk-constrained policies that balance catastrophic unscheduled stoppages against premature overhaul costs~\cite{Aivaliotis_2021, D_He_2025}.
\end{enumerate}

To address these challenges, this paper presents \textbf{KI-SSDM (Knowledge-Informed State-Space Degradation Model)}, a hierarchical state-space framework designed for controller-telemetry prognostics and risk-constrained maintenance decision support. KI-SSDM formalizes domain engineering knowledge into explicit model structures: continuous contact fatigue mechanics serve as the deterministic backbone, while reversible thermal and lubrication variations are captured via a latent Ornstein-Uhlenbeck (OU) nuisance process. By analytically marginalizing the linear-Gaussian nuisance state through Generalised Least Squares (GLS), exact Gaussian telemetry likelihoods are evaluated in $\mathcal{O}(K)$ time via Kalman innovations, while damage and unit frailty are inferred via adaptive quadrature profiling. Structural discrepancies are captured via a deep neural encoder with smooth $C^\infty$ tanh-bounded skip connections ($e^{-b} \le \alpha \le e^b$). A validation-weighted linear pool (stacking) dynamically blends predictive distributions between pure physics and learned corrections, while origin-stratified split-conformal calibration delivers distribution-free finite-sample coverage guarantees. Finally, these calibrated uncertainty bounds feed a downstream engineering decision layer for risk-constrained maintenance scheduling and lifecycle cost optimisation.

\subsection{Related Work and Positioning in Engineering Informatics}
\label{subsec:related}

\paragraph{Physics-Based and Stochastic Degradation Models}
First-principles mechanical models, such as Lundberg-Palmgren fatigue laws~\cite{Lundberg_1947} and Palmgren-Miner damage accumulation rules~\cite{Miner_1945}, provide closed-form interpretability and guarantee monotonic progression~\cite{Lei_2018, H_Li_2024, Yan_2021, N_Zhang_2024, Yucesan_2021}. Integrating empirical degradation curves into dynamic machine representations has also advanced industrial robotics and feed drives~\cite{Aivaliotis_2021, Tai_2023, Wan_2025}. However, they rely on rigid material parameters that vary across batches and cannot adapt to unmodeled lubrication contamination or thermal expansion. While Bayesian particle filters (PFs) track nonlinear degradation under non-Gaussian observations~\cite{Deng_2020, Yan_2021}, their accuracy degrades under observation misspecification, and particle impoverishment frequently destabilizes long-horizon forecasts.

\paragraph{Physics-Informed Machine Learning (PIML) in Engineering Informatics}
Embedding physical constraints into neural architectures has gained substantial traction in engineering informatics and condition monitoring~\cite{Wu_2024, Xu_2023, Yan_2026, Lai_2024}. Contributions in \emph{Advanced Engineering Informatics} have demonstrated the value of physics-informed frameworks, including self-attention PINNs for mechanical remaining life~\cite{Liao_2023} and a two-stage updated digital twin for bearing prognosis~\cite{D_He_2025}; hybrid physics--data models have also been applied to spindles and feed systems~\cite{Tai_2023, Zhang_2025, Yuan_2024}. Other implementations include loss-regularized PINNs with adaptive weighting~\cite{Yan_2026, P_Zhang_2026}, physics-guided autoencoders~\cite{Lai_2024}, and physical information-constrained dynamic wavelets~\cite{Zhao_2026}. While preliminary work explored offline digital-twin modelling and PINN-AIDB regression for ball-screw drives~\cite{PriorWork_MSSP_2026}, that framework relied on heuristic loss penalties without statistical likelihood semantics,  and did not quantify calibrated prediction intervals. In contrast, the present work develops a unified hierarchical state-space formulation with explicit domain knowledge mapping, analytical nuisance marginalisation, validation-weighted distribution pooling, finite-sample conformal uncertainty calibration.

\paragraph{Deep State-Space Models and Conformal UQ for Decision Support}
Deep state-space models (DSSMs) combine explicit latent transition equations with parameterised neural emissions to capture complex temporal dynamics~\cite{Marco_2026, Jiaxin_2025}. Coupling state-space structures with physics-informed learning offers a principled path to model continuous-discrete latent dynamics~\cite{Arnold_2021}. In particular, linear-Gaussian state-space models (LGSSMs) can be trained via Kalman filtering and direct marginal likelihood evaluations~\cite{Marco_2026}. However, existing prognostic DSSMs often rely either on phenomenological latent dynamics lacking mechanical wear backbones or on high-dimensional stochastic latent variables requiring intensive Monte Carlo sampling. In parallel, conformal prediction provides distribution-free, finite-sample coverage under exchangeability~\cite{Vovk_2022, Lei_2018_CP, Angelopoulos_2023, Zhang_CP_2025}. The proposed KI-SSDM leverages origin-stratified split-conformal prediction, calibrating forecast bands across individual operational horizons to ensure finite-sample coverage across fleets, which directly feeds into a risk-constrained decision layer for predictive maintenance.

\subsection{Key Contributions}
\label{subsec:contributions}

The primary scientific and methodological contributions of this work are:
\begin{enumerate}[label=\textbf{C\arabic*.},leftmargin=2.5em]
    \item \textbf{Explicit Knowledge-Informed State-Space Representation}: We establish an explicit mapping from physical mechatronic domain knowledge (ISO~3408-5 contact fatigue~\cite{ISO_3408_5}, dynamic tribology, thermodynamic dissipation, unit frailty, and monotonic irreversibility) into a unified unit-marginal continuous-time state-space architecture, combining monotonic irreversible fatigue exposure $D(t)$ with a stationary, mean-reverting Ornstein-Uhlenbeck nuisance process $B(t)$.
    \item \textbf{Analytic GLS marginalisation and Smooth Bounded Epistemic Discrepancy}: An exact conditional Generalised Least Squares (GLS) reduction integrates $B(t)$ out analytically, yielding closed-form marginal Gaussian telemetry likelihoods evaluated in $\mathcal{O}(K)$ linear time via Kalman innovations recursion, while latent damage $D(t)$ and unit frailty $\eta_u$ are profiled via adaptive quadrature. Structural unmodelled dynamics are captured via a multi-start deep discrepancy encoder whose smooth ($C^\infty$) tanh output confines the log-rate correction to $[-b, +b]$, so that the wear-rate multiplier stays strictly positive and damage accumulation remains monotone ($\mathrm{d}D/\mathrm{d}t \ge 0$) while the likelihood remains differentiable for gradient-based fitting.
    
    \item \textbf{Validation-Weighted Linear-Pool Predictive Stacking}: A validation-optimised mixture of physics-based and discrepancy-based predictive distributions is employed following continuous distribution stacking principles~\cite{Yao_2018}, with weights $(1-w)$ and $w$, respectively. The mixing weight is selected from $w \in \{0,\,0.25,\,0.50,$ $0.75,\,1.00\}$. In the event of tied validation performance, a conservative preference for the physics-based model ($w=0$) is applied to mitigate overfitting.

    \item \textbf{Knowledge-Intensive Decision Support Layer for Maintenance Optimisation}: We extend RUL forecasting into downstream risk-constrained opportunistic maintenance scheduling and lifecycle cost optimisation. By coupling the origin-stratified conformal lower bound $\mathrm{RUL}_{\mathrm{lo}}$ to an industrial cost-risk objective function, the framework kept the empirical failure rate at 3.9\% (planned duty) and reduced expected maintenance cost by 34.7\% (planned) and 31.6\% (probabilistic) relative to periodic $L_{10}$ replacement. The saving comes chiefly from scheduling at a calibrated lower bound: every comparator scheduled the same way saved 30--35\%.
    \item \textbf{Multi-Fleet Benchmark with Multiplicity Control}: Validation across 30 independent synthetic fleets (2,220 simulated units; 480 held-out test units; six degradation scenarios). After Holm adjustment within each duty, KI-SSDM is significantly more accurate than Inverse PINN (planned duty, $p_{\mathrm{Holm}} = 0.0079$) and Deep LSTM (both duties, $p_{\mathrm{Holm}} < 10^{-9}$), and not significantly different from Physics SSM, Bayesian PF ($p_{\mathrm{Holm}} = 0.107$) or Basis Hybrid ($p_{\mathrm{Holm}} = 0.418$). Under load-exponent misspecification (S6) it gains +5.99~h (+29.8\%) over uncorrected physics under planned duty ($p_{\mathrm{Holm}} = 0.0277$ within the S6 family); across S1--S5 the mean gain is +0.02~h.\end{enumerate}

The remainder of this article is organized as follows: Section~\ref{sec:bench_features} describes the industrial BSFD test platform, physics-based digital twin, and feature engineering. Section~\ref{sec:methodology} derives the mathematical formulation of KI-SSDM, including analytic GLS marginalisation, bounded neural discrepancy, and predictive stacking. Section~\ref{sec:conformal_uq} presents the origin-stratified conformal prediction framework. Section~\ref{sec:experiments} details the multi-fleet experimental benchmark. Section~\ref{sec:results} presents benchmark prognostic accuracy, ablation studies, and lifecycle decision support. Section~\ref{sec:discussion} discusses physical interpretability, computational cost, and operational scope. Section~\ref{sec:conclusions} concludes the article. Independent cross-domain validation on the PRONOSTIA bearing benchmark is provided in the accompanying companion Annexure.

% ========================  2. TEST BENCH & FEATURES  ========================
% ========================  2. TEST BENCH & DIGITAL TWIN  ========================
\section{Industrial Test Bench, Digital Twin Modelling, and Feature Engineering}
\label{sec:bench_features}

\subsection{Ball-Screw Feed Drive Test Platform and Edge Setup}
\label{subsec:testbench}

The empirical foundation of this study is an industrial ball-screw feed drive (BSFD) test platform (Fig.~\ref{fig:bsfd_setup}) driven by a Siemens 840Dsl CNC. The mechanical assembly comprises a permanent-magnet synchronous servomotor, a PMI precision ground ball-screw (model R32-10T3-FDIC-990-1152-0.018, grade C3, 10~mm nominal lead, 990~mm active length) with catalogue dynamic load and fatigue ratings~\cite{PMI_Catalog}, a preloaded double-nut carriage, recirculating linear guideways, and a load carriage carrying interchangeable payloads up to 150~kg ($M_{\mathrm{tot}} = 175$~kg including the 25~kg carriage). This configuration represents high-precision feed drives in CNC machining centres and industrial automation~\cite{Deng_2020, Pandhare_2021, Jia_2019}.

\begin{figure*}[pos=!t]
  \centering
  \includegraphics[width=0.88\textwidth, height=0.28\textheight, keepaspectratio]{BSFD_ALT_Testsetup.png}
  \caption{Industrial ball-screw feed drive (BSFD) experimental test bench. Telemetry is streamed via OPC~UA; no external accelerometers are used.}
  \label{fig:bsfd_setup}
\end{figure*}

Telemetry is captured via an embedded OPC~UA server exposing internal drive variables: commanded voltage $u(t)$, motor current $i(t)$, drive-estimated torque $\tau(t) = k_i i(t)$, rotational velocity $\omega(t)$, active power $P(t)$, and table position $x(t)$. Streams are recorded over Gigabit Industrial Ethernet at a sampling cadence of 50--200~ms (median 9.2~Hz).

\subsection{Physics-Based Digital Twin Modelling of the BSFD}
\label{subsec:dt_modeling}

OPC~UA telemetry alone cannot isolate wear-induced friction degradation from dynamic cutting or inertial loads. To decouple these phenomena in real time, a high-fidelity nonlinear digital twin runs on the edge processor, coupling motor electrical dynamics, a two-inertia torsional drivetrain with backlash, linear table mechanics, and dynamic LuGre friction~\cite{Tai_2023, Yuan_2024, Jia_2019, Heydarnia_2025, Wan_2025}.

\paragraph{Motor Electrical Dynamics}
The servomotor electrical subsystem is governed by:
\begin{equation}
    v_m(t) = R_m i(t) + L_m \frac{\mathrm{d}i(t)}{\mathrm{d}t} + k_{\mathrm{emf}} \omega_m(t), \quad \tau_m(t) = k_i i(t),
    \label{eq:elec_dynamics}
\end{equation}
where $v_m(t)$ is armature terminal voltage, $i(t)$ is q-axis current, $R_m$ and $L_m$ are winding resistance and inductance, $k_{\mathrm{emf}}$ is back-EMF constant, $k_i$ is torque sensitivity, and $\omega_m(t)$ is motor angular velocity.

\paragraph{Torsional Drivetrain with Monotone Smooth Backlash Operator}
The flexible coupling connecting the servomotor to the ball screw is modeled as a two-inertia torsional system subject to mechanical backlash. The axial backlash clearance $b_{\mathrm{ax}} = 3.0 \times 10^{-5}$~m translates to angular deadband $\Delta_b = b_{\mathrm{ax}} / i_R \approx 1.886 \times 10^{-2}$~rad, with reduction ratio $i_R = l_s / (2\pi) = 1.5915 \times 10^{-3}$~m/rad for lead $l_s = 10$~mm. To ensure non-negative incremental torsional stiffness across reversals and eliminate unphysical negative-slope instabilities, backlash is represented via a monotone smooth dead-zone operator $\mathcal{P}_{\delta,\varepsilon_b}(\psi)$ based on softplus functions:
\begin{equation}
    \mathcal{P}_{\delta,\varepsilon_b}(\psi) = s_{\varepsilon_b}(\psi - \delta) - s_{\varepsilon_b}(-\psi - \delta),
    \label{eq:backlash_softplus}
\end{equation}
where $s_{\varepsilon_b}(x) = \varepsilon_b \ln\bigl(1 + \exp(x / \varepsilon_b)\bigr)$ with smoothing parameter $\varepsilon_b = 5 \times 10^{-4}$~rad, and $\delta = \Delta_b / 2 \approx 9.43 \times 10^{-3}$~rad denotes angular half-clearance ($b_{\mathrm{ax}} / 2 = 15~\mu\mathrm{m}$). Because $\mathrm{d}\mathcal{P}_{\delta,\varepsilon_b} / \mathrm{d}\psi \ge 0$ everywhere, passive contact dynamics are guaranteed. The coupled drivetrain equations are:

\begin{equation}
\label{eq:tor_motor}
J_m\dot{\omega}_m
= \tau_m-\tau_c-B_m\omega_m.
\end{equation}

\begin{equation}
\label{eq:tor_screw}
\begin{aligned}
\bigl(J_s+m_{\mathrm{tbl}}i_R^2\bigr)
\dot{\omega}_s
&= \tau_c-B_s\omega_s\\
&\quad-i_R\Bigl[
F_{\mathrm{ext}}+F_f(v_{\mathrm{tbl}})\\
&\qquad\quad
+C_a(v_{\mathrm{tbl}})v_{\mathrm{tbl}}
\Bigr].
\end{aligned}
\end{equation}

\begin{equation}
\label{eq:tor_couple}
\begin{aligned}
\tau_c
&= K_t\mathcal{P}_{\delta,\varepsilon_b}
(\theta_m-\theta_s)\\
&\quad+C_t(\omega_m-\omega_s).
\end{aligned}
\end{equation}

where $J_m$ and $J_s$ denote motor and screw polar moments of inertia, $B_m = 1.80 \times 10^{-3}$~N$\cdot$m$\cdot$s/rad and $B_s = 1.20 \times 10^{-3}$~N$\cdot$m$\cdot$s/rad are viscous damping coefficients, $K_t = 1.20 \times 10^4$~N$\cdot$m/rad and $C_t = 2.50$~N$\cdot$m$\cdot$s/rad are coupling stiffness and damping, $\theta_m$ and $\theta_s$ are rotor positions, $m_{\mathrm{tbl}} = 25.0$~kg is unladen carriage mass (reflected table inertia $m_{\mathrm{tbl}} i_R^2 \approx 6.33 \times 10^{-5}$~kg$\cdot$m$^2$, total unladen screw inertia $J_{s,\mathrm{eq}} \approx 2.78 \times 10^{-4}$~kg$\cdot$m$^2$), and $v_{\mathrm{tbl}} = i_R \omega_s$ enforces the rigid screw-nut kinematic constraint. For the laden configuration, $M_{\mathrm{tot}} = 175.0$~kg ($25.0$~kg carriage $+ 150.0$~kg payload), yielding reflected inertia $M_{\mathrm{tot}} i_R^2 \approx 4.43 \times 10^{-4}$~kg$\cdot$m$^2$ and total loaded screw inertia $J_{s,\mathrm{eq},\mathrm{laden}} \approx 6.58 \times 10^{-4}$~kg$\cdot$m$^2$.

\paragraph{Axial Dynamics and Equivalent Modal Coordinates}
In equivalent axial translation coordinates, the unladen feed drive exhibits an effective lumped axial mass of $M_x = J_s / i_R^2 + m_{\mathrm{tbl}} \approx 110.0$~kg ($M_{x,\mathrm{laden}} \approx 260.0$~kg under full $150.0$~kg payload), with carriage translation governed by:
\begin{equation}
    M_x \ddot{x}_{\mathrm{tbl}} = F_m - F_{\mathrm{ext}} - C_a(v_{\mathrm{tbl}}) v_{\mathrm{tbl}} - F_f(v_{\mathrm{tbl}}),
    \label{eq:axial_dynamics}
\end{equation}
where $F_m = \tau_c / i_R$ is axial drive thrust, $F_{\mathrm{ext}}$ is external force, $C_a(v_{\mathrm{tbl}})$ is velocity-dependent damping, and $F_f(v_{\mathrm{tbl}})$ is dynamic friction. In lumped modal coordinates, the reflected load is $\tau_\ell = i_R (M_x \ddot{x}_{\mathrm{tbl}} + C_x \dot{x}_{\mathrm{tbl}} + K_x x_{\mathrm{tbl}} + f_{\mathrm{nl}} + F_{\mathrm{ext}})$~\cite{PriorWork_MSSP_2026}, where $C_x$ and $K_x$ denote translational damping and stiffness, and $f_{\mathrm{nl}}$ aggregates backlash and frictional hysteresis.

\paragraph{Dynamic LuGre Friction Law}
Friction at the ball screw--nut interface and linear guideways is captured via an axial LuGre dynamic friction model~\cite{CanudasDeWit_1995, PriorWork_MSSP_2026, Wan_2025}. The microscopic elastic bristle deflection state $z(t)$ evolves according to:
\begin{equation}
    \dot{z} = v_{\mathrm{tbl}} - \frac{\sigma_0 |v_{\mathrm{tbl}}|}{g(v_{\mathrm{tbl}})} z,
    \label{eq:lugre_state_eq}
\end{equation}
where the velocity-dependent Stribeck function $g(v_{\mathrm{tbl}})$ governs steady-state friction transitions:
\begin{equation}
    g(v_{\mathrm{tbl}}) = F_c + (F_s - F_c) \exp\left[-\left(\frac{v_{\mathrm{tbl}}}{v_s}\right)^2\right],
    \label{eq:stribeck_func}
\end{equation}
with Coulomb friction level $F_c = 12.0$~N, static breakaway friction $F_s = 16.0$~N, and Stribeck velocity $v_s = 0.03$~m/s. The total dynamic friction force $F_f(v_{\mathrm{tbl}})$ and nonlinear axial damping $C_a(v_{\mathrm{tbl}})$ are:
\begin{align}
    F_f(v_{\mathrm{tbl}}) &= \sigma_0 z + \sigma_1 \dot{z} + \sigma_2 v_{\mathrm{tbl}}, \label{eq:lugre_force_eq} \\
    C_a(v_{\mathrm{tbl}}) &= \max\bigl(c_{a0} + c_{a1}|v_{\mathrm{tbl}}|, \; c_{\mathrm{floor}}\bigr), \label{eq:viscous_floor}
\end{align}
where $\sigma_0 = 1.0 \times 10^5$~N/m is bristle contact stiffness, $\sigma_1 = 50.0$~N$\cdot$s/m is bristle damping, $\sigma_2 = 18.5$~N$\cdot$s/m is viscous friction, $c_{a0} = 48.0$~N$\cdot$s/m is base axial damping, $c_{a1} = 15.0$~N$\cdot$s$^2$/m$^2$ is damping slope, and $c_{\mathrm{floor}} = 2.0$~N$\cdot$s/m is a lower bound. Under steady sliding ($\dot{z} = 0$), deflection reaches $z_{\mathrm{ss}} = \frac{g(v_{\mathrm{tbl}})}{\sigma_0} \operatorname{sgn}(v_{\mathrm{tbl}})$, recovering the Stribeck curve: $F_s = 16.0$~N at breakaway and $F_c = 12.0$~N at moderate sliding. In normalized coordinates~\cite{PriorWork_MSSP_2026}, defining $z_a = z / z_{\mathrm{max}} \in [-1, 1]$ with $z_{\mathrm{max}} = F_c / \sigma_0 = 1.2 \times 10^{-4}$~m yields equivalent coefficient $\sigma_0^{\mathrm{norm}} = 12.0$~N, while macroscopic nut stiffness follows catalogue ratings ($K_N = 6.57 \times 10^8$~N/m for R32-10T3-FDIC~\cite{PMI_Catalog}). Differential equations are integrated using a 4th-order Runge-Kutta (RK4) solver with internal sub-stepping of $h = 0.2$~ms (50 sub-steps per 10~ms control cycle).

\paragraph{Cascaded Servo Control Subsystem}
The digital twin replicates the Siemens 840Dsl cascaded control loop. An outer proportional-integral-derivative position controller generates the velocity setpoint $\omega_{\mathrm{ref}}$, while an inner speed controller with LuGre friction feedforward and anti-windup generates the drive reference:
\begin{equation}
    u_{\mathrm{ref}} = \mathrm{sat}_{V_{\max}}\Bigl(K_{p,\omega} e_\omega + K_{i,\omega} \int e_\omega^*\,\mathrm{d}t + K_{\mathrm{ff}} \hat{F}_f\Bigr),
    \label{eq:servo_control}
\end{equation}
where $e_\omega = \omega_{\mathrm{ref}} - \omega_m$, and $\hat{F}_f$ is the digital-twin friction feedforward compensation.

Table~\ref{tab:bsfd_dt_params} summarizes the complete set of identified physical, electrical, and control parameters for the BSFD digital twin.

\begin{table*}[pos=!t]
\centering
\caption{Nonlinear BSFD digital-twin physical, electromechanical, friction, and control parameters~\cite{PriorWork_MSSP_2026}.}
\label{tab:bsfd_dt_params}
\renewcommand{\arraystretch}{1.15}
\setlength{\tabcolsep}{4pt}
\small
\begin{tabularx}{\textwidth}{@{}
  >{\raggedright\arraybackslash}X
  >{\centering\arraybackslash}p{0.08\textwidth}
  >{\centering\arraybackslash}p{0.18\textwidth}
  >{\raggedright\arraybackslash}X
  >{\centering\arraybackslash}p{0.08\textwidth}
  >{\centering\arraybackslash}p{0.18\textwidth}
@{}}
  \toprule
  \multicolumn{3}{@{}l}{\textbf{Kinematics, Inertia \& Torsional Coupling}} &
  \multicolumn{3}{l@{}}{\textbf{Motor--Electrical Subsystem}}\\
  \midrule
  Parameter Specification & Symbol & Value (unit) &
  Parameter Specification & Symbol & Value (unit) \\
  \midrule
  Lead / reduction ratio & $i_R$ & $1.5915{\times}10^{-3}$~m/rad &
  Armature winding resistance & $R_m$ & $1.0~\Omega$ \\
  Unladen carriage mass & $m_{\mathrm{tbl}}$ & $25.0$~kg &
  Armature winding inductance & $L_m$ & $5.50{\times}10^{-3}$~H \\
  Laden table mass & $M_{\mathrm{tot}}$ & $175.0$~kg &
  Torque constant sensitivity & $k_i$ & $0.82$~N$\cdot$m/A \\
  Motor rotor polar inertia & $J_m$ & $7.70{\times}10^{-4}$~kg$\cdot$m$^{2}$ &
  Back-EMF constant & $k_{\mathrm{emf}}$ & $0.15$~V$\cdot$s/rad \\
  Screw rotor polar inertia & $J_s$ & $2.15{\times}10^{-4}$~kg$\cdot$m$^{2}$ &
  Drive voltage clamp & $V_{\max}$ & $400$~V \\
  Torsional coupling stiffness & $K_t$ & $1.20{\times}10^4$~N$\cdot$m/rad &
  Rotational speed clamp & $\Omega_{\max}$ & $750$~rad/s \\
  Torsional coupling damping & $C_t$ & $2.50$~N$\cdot$m$\cdot$s/rad &
  Velocity setpoint clamp & $V_{\mathrm{ref,max}}$ & $1.2$~m/s \\
  Motor / screw damping & $B_m,\allowbreak B_s$ & $1.80{\times}10^{-3},\allowbreak 1.20{\times}10^{-3}$~N$\cdot$m$\cdot$s/rad &
  Actuator time constant & $\tau_u$ & $0.01$~s \\
  Backlash half-clearance & $\delta$ & $9.43{\times}10^{-3}$~rad ($b_{\mathrm{ax}} = 30~\mu\mathrm{m}$) &
  Numerical solver sub-step & $h$ & $0.2$~ms (RK4) \\
  Backlash smooth factor & $\varepsilon_b$ & $5.00{\times}10^{-4}$~rad &
  Nominal catalogue lead & $l_s$ & $10.0$~mm \\
  \midrule
  \multicolumn{3}{@{}l}{\textbf{Axial Viscous \& Dynamic LuGre Friction}} &
  \multicolumn{3}{l@{}}{\textbf{Cascaded Controller \& Friction Compensation}}\\
  \midrule
  Base viscous coefficient & $c_{a0}$ & $48.0$~N$\cdot$s/m &
  Position PID gains & $K_{p,p},\allowbreak K_{i,p},\allowbreak K_{d,p}$ & $600,\allowbreak 3.0{\times}10^{4},\allowbreak 60$ \\
  Speed-dependent slope & $c_{a1}$ & $15.0$~N$\cdot$s$^2$/m$^2$ &
  Speed PI gains & $K_{p,\omega},\allowbreak K_{i,\omega}$ & $3.0,\allowbreak 120$ \\
  Viscous floor bound & $c_{\mathrm{floor}}$ & $2.0$~N$\cdot$s/m &
  Anti-windup back-calc & $\tau_{\mathrm{aw},p},\allowbreak \tau_{\mathrm{aw},\omega}$ & $10,\allowbreak 20$~s \\
  Coulomb friction level & $F_c$ & $12.0$~N &
  Friction feedforward & $\alpha_{\mathrm{ff}},\allowbreak K_{\mathrm{boost}},\allowbreak v_{\mathrm{ss}}$ & $0.2,\allowbreak 0.8,\allowbreak 0.02$~m/s \\
  Static breakaway friction & $F_s$ & $16.0$~N &
  Nut axial stiffness & $K_N$ & $6.57{\times}10^8$~N/m ($67~\mathrm{kgf/\mu m}$) \\
  Stribeck knee velocity & $v_s$ & $0.03$~m/s &
  Basic dynamic load rating & $C_a$ & $26.5$~kN \\
  Bristle elastic stiffness & $\sigma_0$ & $1.0{\times}10^{5}$~N/m ($\sigma_0^{\mathrm{norm}} = 12.0$~N) &
  Basic static load rating & $C_{0a}$ & $68.7$~kN \\
  Bristle micro-damping & $\sigma_1$ & $50.0$~N$\cdot$s/m &
  Screw nominal diameter & $d_s$ & $32.0$~mm \\
  Extra viscous friction & $\sigma_2$ & $18.5$~N$\cdot$s/m &
  Nut model designation & --- & PMI R32-10T3-FDIC~\cite{PMI_Catalog} \\
  \bottomrule
\end{tabularx}
\end{table*}

\subsection{Digital Twin Qualification and Experimental PRBS Tracking Validation}
\label{subsec:prbs_validation}

To verify that the digital twin provides authentic physics-grounded states rather than idealized approximations, dynamic qualification was executed using a Pseudo-Random Binary Sequence (PRBS) position excitation protocol on the physical Siemens 840Dsl test bench and simulated concurrently on the digital twin. The PRBS sequence switched carriage position between limits $(X_0, X_1) = (0, 0.70)$~m with a bit duration of $T_b = 15$~s, exciting reversal dynamics and quasi-steady velocity plateaus.

\begin{figure}[pos=!t]
  \centering
  \includegraphics[width=0.48\textwidth, keepaspectratio]{PRBS_Pulse_Supp.png}
  
\caption{Experimental PRBS tracking validation over 700~mm traversals. The main panel compares commanded
and digital-twin-predicted positions; the inset (310--325~s) highlights steady-state plateau
tracking ($\pm0.2$~mm). Latency-corrected RMSE is 0.15~mm (0.021\% of stroke), while the
uncorrected RMSE of 1.18~mm reflects OPC~UA communication latency (96~ms).}

  \label{fig:prbs_validation}
\end{figure}

Fig.~\ref{fig:prbs_validation} illustrates the dynamic response over a representative segment. The qualification metrics confirm high closed-loop structural fidelity:
\begin{enumerate}[label=(\alph*),leftmargin=2.2em]
    \item \textbf{Dynamic Position Tracking Accuracy}: Across quasi-steady plateau segments (last 30\% of each bit), the raw measured position RMSE was $1.18$~mm (0.17\% of stroke). Decomposing this residual reveals that $\approx 96$~ms corresponds to the deterministic OPC~UA Ethernet transport latency during dynamic velocity ramps at $0.30$~m/s traversal speed, whereas settled dwell plateaus reflect closed-loop servo stiffness. Applying a 96~ms timestamp alignment across held-out runs reduces tracking RMSE to \textbf{$0.15$~mm (0.021\% of stroke)}~\cite{PriorWork_MSSP_2026}, confirming that the digital twin tracks physical plant kinematics with high fidelity (for comparison, JIS C3 ground lead tolerance is $e_{300} \le 8\text{--}12~\mu\mathrm{m} / 300\mathrm{~mm}$~\cite{PMI_Catalog}; the observed residual of $0.15$~mm over the full 700~mm stroke reflects uncalibrated structural compliance, mounting runout, and pre-sliding bristle deflection).
    \item \textbf{Electromechanical Feature Correlation}: Across the full PRBS trajectory, motor current achieves a Pearson correlation of \textbf{$\rho = 0.9991$} ($\mathrm{RMSE} = 0.31$~A), rotational velocity achieves \textbf{$\rho = 0.9986$} ($\mathrm{RMSE} = 1.42$~rad/s), and torque innovation error $e_k = \tau - h_\psi(\mathbf{r}_k)$ exhibits an RMSE of $0.12$~$\mathrm{N\cdot m}$ ($\sigma_e = 0.12$~$\mathrm{N\cdot m}$) with baseline bias $\bar{e}_k = 0.08$~$\mathrm{N\cdot m}$ (representing pristine tare friction). Correlation $\rho$ is omitted for $e_k$ because, as an orthogonal innovation sequence, it behaves as a zero-mean process orthogonal to the reference trajectory.
\end{enumerate}
Table~\ref{tab:prbs_metrics} presents the quantitative dynamic qualification metrics across $n = 605$ steady-state traversal windows evaluated across 10 repeated experimental trials.


%========================================================
% PRBS TABLE - SINGLE COLUMN FORMAT
% All column widths are independently adjustable
%========================================================
\begin{table}[!t]
\centering

\caption{Quantitative dynamic PRBS tracking
qualification metrics across $n=605$ steady-state
traversal windows (quasi-steady plateau: final
30\% of each bit).}

\label{tab:prbs_metrics}

%========================================================
% USER CONFIGURABLE COLUMN WIDTHS
%========================================================
\newlength{\prbsColA}
\newlength{\prbsColB}
\newlength{\prbsColC}
\newlength{\prbsColD}
\newlength{\prbsColE}

\setlength{\prbsColA}{3.00cm} % Signal / Metric
\setlength{\prbsColB}{1.00cm} % RMSE
\setlength{\prbsColC}{1.00cm} % MAE
\setlength{\prbsColD}{1.00cm} % Max Error
\setlength{\prbsColE}{1.00cm} % Pearson rho

%========================================================
% TABLE FORMATTING
%========================================================
\renewcommand{\arraystretch}{1.15}
\setlength{\tabcolsep}{2pt}

\footnotesize

%========================================================
% AUTOMATIC SINGLE-COLUMN WIDTH CONSTRAINT
%========================================================
\begin{adjustbox}{max width=\columnwidth}

\begin{tabular}{@{}
>{\raggedright\arraybackslash}p{\prbsColA}
>{\centering\arraybackslash}p{\prbsColB}
>{\centering\arraybackslash}p{\prbsColC}
>{\centering\arraybackslash}p{\prbsColD}
>{\centering\arraybackslash}p{\prbsColE}
@{}}

\toprule

\textbf{Signal / Metric}
& \textbf{RMSE}
& \textbf{MAE}
& \textbf{Max Error}
& \textbf{Pearson} $\bm{\rho}$ \\

\midrule

\multicolumn{5}{@{}l}{%
\textit{Level 1 Digital Twin Feature Signals}} \\

\addlinespace[2pt]

Motor current $i$ (A)
& 0.31 & 0.24 & 0.92 & 0.9991 \\

Rotational velocity
$\omega_m$ (rad/s)
& 1.42 & 1.07 & 4.10 & 0.9986 \\

Torque innovation
$e_k$ (N$\cdot$m)$^\dagger$
& 0.12 & 0.08 & 0.35 & --- \\

\midrule

\multicolumn{5}{@{}l}{%
\textit{Table Kinematics (Verification Reference)}} \\

\addlinespace[2pt]

$x_{\mathrm{tbl}}$ raw
OPC UA (mm)$^\ddagger$
& 1.18 & 0.91 & 4.24 & 0.9998 \\

$x_{\mathrm{tbl}}$
latency-corrected (mm)$^*$
& $\mathbf{0.15}$
& ---
& ---
& $\approx 1.000$ \\

\bottomrule

\end{tabular}
\end{adjustbox}

%========================================================
% FOOTNOTES - CONFINED TO SINGLE COLUMN WIDTH
%========================================================
\vspace{3pt}

\begin{minipage}{\columnwidth}
\scriptsize
\raggedright

$^\dagger$ Pristine baseline torque innovation bias
is $\bar{e}_k=0.08~\mathrm{N\cdot m}$ ($\sigma_e = 0.12~\mathrm{N\cdot m}$);
as an orthogonal innovation process, correlation with the setpoint trajectory is omitted.
Degradation is detected by sustained monotonic departure from this baseline.

\par\vspace{2pt}

$^\ddagger$ Raw position RMSE includes
deterministic $\approx96$~ms Ethernet
transport latency at $0.30$~m/s
traversal velocity ($0.30 \times 0.096 \approx 28.8$~mm peak dynamic ramp lag), plus
$\pm0.2$~mm pre-sliding micro-deflection.

\par\vspace{2pt}

$^*$ Latency-corrected RMSE of
$0.15$~mm ($0.021\%$ stroke)
validates dynamic closed-loop servo
fidelity. For reference, JIS C3 ground
lead tolerance is $e_{300} \le 8\text{--}12~\mu\mathrm{m}/300~\mathrm{mm}$,
while standard transport-grade screws span $e_p \le 52~\mu\mathrm{m}/300~\mathrm{mm}$.

\end{minipage}

\end{table}


\paragraph{Hierarchical Distinction: Electromechanical Digital Twin vs.\ State-Space Degradation Model}
To prevent methodological conflation, we distinguish between two hierarchical modelling levels in our architecture:
\begin{enumerate}[label=(\roman*),leftmargin=2.2em]
    \item \textbf{High-Frequency Electromechanical Digital Twin (Level 1)}: A continuous-time electromechanical ODE system (Eqs.~\eqref{eq:elec_dynamics}--\eqref{eq:servo_control}) integrated with explicit RK4 on a 100~Hz grid with $h = 0.2$~ms sub-steps. Its plant states $\mathbf{x}^{\mathrm{DT}}$ are the q-axis current $i_q$, rotor and screw coordinates $(\theta_m,\omega_m,\theta_s,\omega_s)$ and the dynamic LuGre bristle deflection $z$ ($\tau_{\mathrm{bristle}} = \sigma_1/\sigma_0 = 0.5$~ms under $\sigma_0 = 1.0 \times 10^5$~N/m); the cascaded controller adds its integrator and actuator-lag states. From these it outputs, at each window $k$, the 6-dimensional feature vector $\mathbf{Z}_k$ of Eq.~\eqref{eq:dt_state_vec}. The PRBS validation confirms plant replication with $>99.8\%$ correlation before its parameters are permanently frozen.
    \item \textbf{Macroscopic State-Space Degradation Model KI-SSDM (Level 2)}: A continuous-time stochastic state-space model operating on macroscopic operational timescales (hours to thousands of hours). It tracks irreversible mechanical fatigue damage $D(t) \in [0, 1]$ alongside a latent reversible Ornstein-Uhlenbeck nuisance process $B(t)$ representing thermal/vibrational variations. Ingesting the 23-dimensional feature vector $\bPhi_k$ from Level~1, it analytically marginalises $B(t)$ via GLS and evaluates conformal prediction bounds for RUL.
\end{enumerate}
Once qualified, all digital-twin parameters were permanently frozen, ensuring downstream features reflect authentic physical states without empirical drift.

\subsection{Real-Time Telemetry Resampling, Windowing, and Feature Engineering}
\label{subsec:realtime_features}

The real-time edge processing pipeline ingests controller telemetry asynchronously and synchronizes it for recursive state-space inference through four sequential operations:

\paragraph{Cadence Normalization and Sliding Windowing}
In industrial CNC machine tool practice, real-time OPC~UA servers deliver telemetry at an asynchronous cadence of approximately $100$~ms ($\approx 10$~samples per second, observed median $9.2$~Hz on the Siemens 840Dsl). Because dynamic friction and electromechanical integration require higher temporal resolution, this raw telemetry stream is linearly upsampled onto a standardized $100$~Hz uniform temporal grid ($100$~samples per second, $10$~ms spacing). Sliding analysis windows of duration $T_w = 1.0$~s ($N_w = 100$ samples) with 50\% overlap ($\rho = 0.50$) define the fundamental recursive processing timestep:
\begin{equation}
    \Delta t = \frac{(1 - \rho) T_w}{3600} \approx 1.389 \times 10^{-4}~\mathrm{h} \quad (500~\mathrm{ms~hop}).
    \label{eq:dt_step}
\end{equation}

\paragraph{Direction-Reversal Flagging}
Because the Siemens 840Dsl streams unsigned motor velocity, kinematic reversals are detected directly from discrete table position increments:
\begin{equation}
    \hat{n}_{\mathrm{rev},k} = \sum_{n=1}^{N_w-2} \mathbf{1}\bigl[\mathrm{sgn}(\Delta x_{n+1}) \ne \mathrm{sgn}(\Delta x_n)\bigr],
    \label{eq:reversal_flag}
\end{equation}
with zero increments assigned the preceding non-zero sign. Reversal counts parameterise kinematic duty factors $f_{\mathrm{op}}(t)$ in the fatigue backbone and are tracked continuously without dropping traversal windows.

\paragraph{Digital Twin Parallel State Generation}
At each window $k$, the digital twin evaluates voltage setpoint $u_k$ to produce the 6-dimensional simulated physical state vector:
\begin{equation}
    \mathbf{Z}_k = \bigl[\omega_k^Z, \; \dot{\omega}_k^Z, \; \tau_{\mathrm{fric},k}^Z, \; \tau_{m,k}^Z, \; i_k^Z, \; u_k^Z\bigr]^\top \in \R^6,
    \label{eq:dt_state_vec}
\end{equation}
supplying dynamic friction torque $\tau_{\mathrm{fric},k}^Z$ and angular acceleration $\dot{\omega}_k^Z$ that cannot be sensed directly from raw controller telemetry.

\paragraph{Disturbance Decoupling and Bagged Innovation Energy}
Total measured motor torque $\tau_k$ is decoupled into external mechanical disturbance torque $\tilde{\tau}_{\mathrm{ext},k}$ by subtracting simulated inertial, viscous, and friction components:
\begin{equation}
    \tilde{\tau}_{\mathrm{ext},k} = \tau_k - \bigl(J_{\mathrm{tot}} \dot{\omega}_k^Z + i_R \hat{F}_{f,k} + B_{\mathrm{tot}} \omega_k^Z\bigr),
    \label{eq:dist_torque}
\end{equation}
where $J_{\mathrm{tot}} = J_m + J_s + m_{\mathrm{tbl}} i_R^2$ is total drivetrain inertia and $B_{\mathrm{tot}} = B_m + B_s$. To detect microscopic wear emergence, a 200-tree bagged regressor $h_\psi(\mathbf{r}_k)$, trained strictly on healthy baseline traversals of training units, predicts baseline torque from transformed twin states $\mathbf{r}_k = [1, \log(|\dot{\omega}_k^Z|+\varepsilon), \log(|\tau_{\mathrm{fric},k}^Z|+\varepsilon), \dots]^\top$. The causal torque innovation residual $e_k = \tau_k^{\mathrm{meas}} - h_\psi(\mathbf{r}_k)$ and its moving-average energy $E_k = \mathrm{MA}_{\mathrm{WMA}}\{e_k^2\}$ serve as sensitive indicators of raceway fatigue initiation.

Assembling all telemetry channels, digital-twin outputs, and residual innovations yields the 23-dimensional prognostic observation vector $\bPhi_k \in \R^{23}$. To satisfy the foundational requirement of engineering informatics for explicit domain knowledge representation~\cite{Liao_2023, D_He_2025}, Table~\ref{tab:phi_blocks} classifies each mechatronic domain knowledge component according to the informed machine learning taxonomy of von Rueden et al.~\cite{vonRueden_2023} across knowledge source, formal mathematical representation, integration point, and epistemic role (Part~A), alongside the diagnostic taxonomy of the 23-channel observation vector (Part~B). In our state-space framework, these 23 features fulfill distinct roles: (1) \emph{Primary emission channel}: the torque innovation residual $e_k$ is fitted by the linear observation model in Eq.~\eqref{eq:observation_model}; (2) \emph{Discrepancy covariates}: feed velocity $|v_{\mathrm{tbl}}|$, torque load $\tau_m$, and kinematic acceleration serve as inputs $\mathbf{u}(t)$ to the bounded neural discrepancy network; and (3) \emph{Diagnostic indicators}: statistical dispersion and spectral power features serve as supervisory health and regime-classification inputs.

\begin{table*}[pos=!t]
\centering
\caption{Knowledge representation and feature architecture in KI-SSDM: (a) Domain knowledge mapping across the taxonomy of von Rueden et al.~\cite{vonRueden_2023}; (b) Structure and diagnostic roles of the 23-dimensional feature vector $\bPhi_k$.}
\label{tab:phi_blocks}
\renewcommand{\arraystretch}{1.15}
\setlength{\tabcolsep}{4pt}
\small
\begin{tabularx}{\textwidth}{@{}p{3.2cm} p{4.4cm} p{4.2cm} X@{}}
\toprule
\multicolumn{4}{@{}l}{\textbf{(a) Explicit Epistemic Knowledge Representation Mapping (Taxonomy of von Rueden et al.~\cite{vonRueden_2023})}} \\
\midrule
\textbf{Knowledge Concept \& Source} & \textbf{Formal Representation} & \textbf{Integration Point \& Architecture} & \textbf{Epistemic \& Decision Role} \\
\midrule
Contact Fatigue Law \newline (ISO~3408-5~\cite{ISO_3408_5, PMI_Catalog}; Hertzian mechanics) & Deterministic ODE: $\frac{\mathrm{d}D_0}{\mathrm{d}t} = c_u \left(\frac{F}{F_{\mathrm{ref}}}\right)^p \frac{3600|\omega_m|}{2\pi} f_{\mathrm{op}}$ & Hypothesis set: Latent damage transition backbone & Anchors long-horizon wear progression; prevents unphysical extrapolation. \\
Wear Irreversibility \newline (2nd Law of Thermodynamics; non-healing wear) & Monotonicity inequality: $\frac{\mathrm{d}D}{\mathrm{d}t} \ge 0$; positive increments $\Delta D_k \ge 0$ & Hypothesis set: Continuous-time positive drift constraint & Enforces thermodynamic irreversibility; eliminates spontaneous damage healing. \\
Reversible Nuisance \newline (Thermal expansion, viscous grease warming $t_{1/2} \approx 12.5$~h) & Mean-reverting Ornstein-Uhlenbeck SDE: $\mathrm{d}B(t) = -\theta_B B(t)\mathrm{d}t + \sigma_B \mathrm{d}W_B$ & Latent state space: Analytically marginalized via $\mathcal{O}(K)$ GLS & Decouples reversible thermal/operational fluctuations from cumulative wear. \\
Unit-to-Unit Frailty \newline (Metallurgical variance, preload \& tolerance dispersion) & Hierarchical log-normal prior: $\log c_u = \mu_C + \eta_u$, $\eta_u \sim \mathcal{N}(0, \sigma_C^2)$ & Prior distribution: Gauss-Hermite quadrature profiling & Regularizes fleet dispersion; shares statistical strength across manufacturing batches. \\
Bounded Authority \newline (Epistemic uncertainty under novel operating regimes) & Smooth bounded mapping: $\alpha(t) = \exp(\rho)$, with $\rho = b \tanh(\dots) \in [-b, b]$ & Hypothesis set: $C^\infty$ smooth tanh-gated multiplicative skip network & Restricts neural correction to $[e^{-b}, e^b] = [0.74, 1.35]$ ($b=0.30$); guarantees non-divergence. \\
Model-Form Risk Aversion \newline (Overfitting risk when nominal physics is well-specified) & Validation linear pool: $p = (1-w)p_{\mathrm{phys}} + w p_{\mathrm{disc}}$ (tie-break $w=0$) & Decision \& post-processing layer: validation log-RUL loss; exact ties $\to$ physics; 1-SE gate for $\kappa$ & Zero-harm fallback: defaults to pure physics if neural corrections do not improve the validation loss. \\
\midrule
\multicolumn{4}{@{}l}{\textbf{(b) 23-Dimensional Prognostic Feature Observation Vector $\bPhi_k$}} \\
\midrule
\textbf{Feature Block} & \textbf{Data Source} & \textbf{Dim} & \textbf{Included Features \& Diagnostic Role} \\
\midrule
Statistical & OPC UA Telemetry & 7 & $\sigma_\tau, \tau_{\mathrm{RMS}}, \sigma_i, i_{\mathrm{RMS}}, \sigma_\omega, \sigma_P, \mathrm{ptp}(P)$ --- Torque/power dispersion and RMS reflecting raceway micro-roughness. \\
Kinematic & OPC UA Telemetry & 6 & $\bar{u}, \bar{i}, \bar{\tau}, \bar{P}, \bar{\omega}, \bar{x}$ --- Traversal means characterising axis operational setpoints and kinematic regimes. \\
Digital Twin & Dynamic Simulation & 6 & $\omega^Z, \dot{\omega}^Z, \tau_{\mathrm{fric}}^Z, \tau_m^Z, i^Z, u^Z$ --- Baseline LuGre friction, speed, and torque from healthy nominal twin. \\
Residuals & Hybrid Reconstruction & 4 & $\tilde{\tau}_{\mathrm{ext}}, \hat{\tau}, e, E$ --- Decoupled disturbance torque, bagged baseline torque, innovation residual $e$, and energy $E$. \\
\bottomrule
\end{tabularx}
\end{table*}

\begin{figure*}[pos=!t]
  \centering
  \includegraphics[width=0.98\textwidth, keepaspectratio]{PI_DSSDM_Framework.png}
  \caption{End-to-end KI-SSDM workflow: (A) OPC~UA telemetry ingestion; (B) digital-twin feature extraction ($\bPhi_k \in \R^{23}$); (C) state-space decomposition into wear $D(t)$ and OU nuisance $B(t)$; (D) analytic GLS marginalisation; (E) bounded discrepancy learning; (F) predictive stacking; and (G) origin-stratified conformal calibration.}
  \label{fig:framework}
\end{figure*}


\section{Mathematical Formulation of KI-SSDM}
\label{sec:methodology}

\subsection{Physical Fatigue Damage Backbone}
\label{subsec:physics_backbone}

Degradation in rolling-element feed drives under variable axial loads and rotational velocities is governed by continuous contact-stress accumulation. Let $D(t) \in [0, \Dcrit]$ denote the dimensionless cumulative fatigue exposure of a machine unit, normalized such that $D(0) = 0$ corresponds to a pristine component and $\Dcrit = 1.0$ represents fatigue spalling or structural pitting failure. Under classical Lundberg-Palmgren rolling contact fatigue (RCF) mechanics~\cite{Lundberg_1947}, Palmgren--Miner linear damage accumulation~\cite{Miner_1945} and ISO~3408-5 dynamic load formulations~\cite{ISO_3408_5, PMI_Catalog}, the continuous deterministic rate of damage accumulation per operating hour $t_h$ is:
\begin{equation}
    \frac{\mathrm{d}D_0(t_h)}{\mathrm{d}t_h} = c_u \left(\frac{F(t_h)}{F_{\mathrm{ref}}}\right)^p \cdot \frac{3600 |\omega_m(t_h)|}{2\pi} \cdot f_{\mathrm{op}}(t_h),
    \label{eq:miner_damage}
\end{equation}
where $F(t_h)$ is the dynamic axial contact load magnitude, $F_{\mathrm{ref}}$ is a reference load normalization anchor (ensuring dimensional consistency when the exponent $p$ varies), $p$ is the contact fatigue exponent (nominally $p=3.0$ for ball bearings and ball screws under ideal elastohydrodynamic lubrication per ISO 3408-5~\cite{ISO_3408_5}), $|\omega_m(t_h)|$ is the absolute motor rotational speed in rad/s (so that $3600|\omega_m|/(2\pi)$ is the number of revolutions per hour), $c_u$ is the unit-specific fatigue damage coefficient per revolution at reference load, and $f_{\mathrm{op}}(t_h) \ge 1$ encapsulates kinematic reversing factors. The use of unsigned rotational speed $|\omega_m|$ guarantees that damage accumulates monotonically ($\mathrm{d}D_0 / \mathrm{d}t_h \ge 0$) regardless of travel direction.

To account for unit-to-unit manufacturing tolerances, metallurgical inclusions, and assembly variations across an industrial fleet, $c_u$ is parameterised by a log-normal distribution:
\begin{equation}
    \log c_u = \mu_C + \eta_u, \quad \eta_u \sim \mathcal{N}(0, \sigma_C^2), \quad u = 1, \dots, N_{\mathrm{units}},
    \label{eq:unit_frailty}
\end{equation}
where $\mu_C$ denotes the population mean of $\log c_u$, $\sigma_C^2$ is the frailty variance, and $\eta_u$ is a mean-zero unit random effect. Under this log-normal specification, the population arithmetic mean is $\mathbb{E}[c_u] = \exp(\mu_C + \sigma_C^2 / 2)$.

\subsection{Latent Reversible Nuisance Dynamics}
\label{subsec:nuisance_b}

Sensor telemetry acquired from industrial controllers (OPC~UA streams capturing motor torque, q-axis current, and drive power) fluctuates due to friction temperature shifts, speed reversals, and transient cutting loads. These variations do not reflect irreversible material loss. KI-SSDM decouples these effects by introducing a zero-mean, stationary Ornstein-Uhlenbeck (OU) nuisance process $B(t)$, providing a continuous-time state-space parameterisation of reversible disturbances distinct from monotonic degradation~\cite{Zhang_2023_RESS}:
\begin{equation}
    \mathrm{d}B(t) = -\theta_B B(t)\,\mathrm{d}t + \sigma_B\,\mathrm{d}W_B(t),
    \label{eq:ou_process}
\end{equation}
where $\theta_B > 0$ is the mean-reversion rate (with mean-reversion half-life $t_{1/2} = \ln(2) / \theta_B$, e.g.\ $12.5$~h for the benchmark value $\theta_B = 1/18\,\mathrm{h}^{-1}$, operating on thermal-reversal timescales), $\sigma_B > 0$ is the diffusion intensity, and $W_B(t)$ is a standard Brownian motion independent of structural wear. For any discrete sampling interval $\Delta t$, the discrete transition density of $B_k = B(t_k)$ is Gaussian:
\begin{equation}
    B_{k+1} \mid B_k \sim \mathcal{N}\Bigl(B_k e^{-\theta_B \Delta t}, \; \frac{\sigma_B^2}{2\theta_B}\bigl(1 - e^{-2\theta_B \Delta t}\bigr)\Bigr).
    \label{eq:ou_transition}
\end{equation}
To ensure structural identifiability and prevent trade-offs between nuisance amplitude and observation scale, we standardise $B$ to unit stationary variance, $B(0) \sim \mathcal{N}(0, 1)$, so that its amplitude is carried by the channel loadings $\gamma_j$. Equivalently $\sigma_B^2 = 2\theta_B$, and the transition variance in Eq.~\eqref{eq:ou_transition} reduces to $1 - e^{-2\theta_B \Delta t}$.

\subsection{Exact Analytic GLS marginalisation and Kalman Innovations}
\label{subsec:gls_marginalization}

The observation model links $J = 3$ controller telemetry channels (q-axis current, speed, and the causal torque innovation residual $e_k = \tau_k^{\mathrm{meas}} - h_\psi(\mathbf{r}_k)$, which isolates progressive tribological degradation from dynamic electromechanical inertial loads) to the shared latent states:
\begin{equation}
\begin{aligned}
    y_j(t_k) &= \boldsymbol{\beta}_j^\top \mathbf{x}_k + \gamma_j B(t_k) + \epsilon_{kj},\\
    \epsilon_{kj} &\sim \mathcal{N}(0, \sigma_j^2), \qquad j = 1, \dots, J,
\end{aligned}
    \label{eq:observation_model}
\end{equation}
where $\mathbf{x}_k = [1, D(t_k), D(t_k)^2, \log F_k, \log v_k, T_k]^\top$ collects the damage terms and the operating covariates (log axial load, log feed speed and temperature), $\boldsymbol{\beta}_j$ are the emission coefficients of channel $j$ (the quadratic coefficient is ridge-penalised), $\gamma_j \in [-2, 2]$ are channel loadings on $B$, and $\sigma_j$ are channel noise standard deviations. Writing $\mathbf{y}_k = [y_1(t_k), \dots, y_J(t_k)]^\top$, $\boldsymbol{\gamma} = [\gamma_1, \dots, \gamma_J]^\top$, $\boldsymbol{\Sigma}_\epsilon = \mathrm{diag}(\sigma_1^2, \dots, \sigma_J^2)$ and collecting all emission coefficients in $\boldsymbol{\beta} \in \R^{6J}$, the conditional mean is $\mathbf{X}_k(D)\boldsymbol{\beta}$ with $\mathbf{X}_k(D) = \mathbf{I}_J \otimes \mathbf{x}_k^\top$. The remaining telemetry signals ($P, x$) and digital-twin states ($\mathbf{Z}_k$) enter the 23-dimensional feature vector $\bPhi_k$ used for neural discrepancy evaluation and supervisory diagnostics.

Because $B(t)$ is linear-Gaussian with unit stationary variance, the trajectory $\mathbf{B}_u = [B(t_{u,1}), \dots, B(t_{u,K_u})]^\top$ of unit $u$ has the exact joint covariance matrix $\boldsymbol{\Sigma}_{B,u} \in \R^{K_u \times K_u}$ with elements
\begin{equation}
    \bigl[\boldsymbol{\Sigma}_{B,u}\bigr]_{k,l} = \exp\bigl(-\theta_B |t_{u,k} - t_{u,l}|\bigr).
    \label{eq:sigma_b}
\end{equation}

Conditioning on the latent damage trajectory $\mathbf{D}_u$ and fixed parameters, the stacked observation vector $\mathbf{y}_u = [\mathbf{y}_1^\top, \dots, \mathbf{y}_{K_u}^\top]^\top \in \R^{K_u J}$ follows an exact multivariate Gaussian distribution:
\begin{equation}
\begin{aligned}
    \mathbf{y}_u \mid \mathbf{D}_u &\sim \mathcal{N}\bigl(\mathbf{X}_u \boldsymbol{\beta}, \; \boldsymbol{\Sigma}_{y,u}\bigr),\\
    \boldsymbol{\Sigma}_{y,u} &= \boldsymbol{\Sigma}_{B,u} \otimes \boldsymbol{\gamma}\boldsymbol{\gamma}^\top + \mathbf{I}_{K_u} \otimes \boldsymbol{\Sigma}_\epsilon,
\end{aligned}
    \label{eq:marginal_gaussian}
\end{equation}
where $\mathbf{X}_u$ stacks $\mathbf{X}_k(D)$ over $k = 1, \dots, K_u$.

When emission coefficients $\boldsymbol{\beta}$ are shared across an industrial training fleet of $U$ units, the pooled Generalised Least Squares (GLS) estimator is resolved analytically:
\begin{equation}
    \hat{\boldsymbol{\beta}} = \left(\sum_{u=1}^U \mathbf{X}_u^\top \boldsymbol{\Sigma}_{y,u}^{-1} \mathbf{X}_u\right)^{-1} \left(\sum_{u=1}^U \mathbf{X}_u^\top \boldsymbol{\Sigma}_{y,u}^{-1} \mathbf{y}_u\right).
    \label{eq:gls_solution}
\end{equation}

Explicitly factorising the dense $K_u J \times K_u J$ covariance matrix $\boldsymbol{\Sigma}_{y,u}$ costs $O(K_u^3)$ and is impractical for long records. Because the OU nuisance process defines a linear-Gaussian state-space model with a scalar state, the marginal likelihood is evaluated exactly in $O(K_u)$ time via the Kalman innovations representation~\cite{Durbin_2012}:
\begin{equation}
\begin{aligned}
    \log p(\mathbf{y}_u \mid \mathbf{D}_u) = -\frac{1}{2}\sum_{k=1}^{K_u} \Bigl(&J\log(2\pi) + \log\lvert\mathbf{F}_k\rvert\\
    &+ \mathbf{v}_k^\top \mathbf{F}_k^{-1} \mathbf{v}_k\Bigr),
\end{aligned}
    \label{eq:kalman_innovations}
\end{equation}
where $\mathbf{v}_k = \mathbf{y}_{u,k} - \hat{\mathbf{y}}_{k|k-1} \in \R^J$ is the pre-update innovation and $\mathbf{F}_k = \boldsymbol{\gamma}\boldsymbol{\gamma}^\top P_{k|k-1} + \boldsymbol{\Sigma}_\epsilon$ is its covariance, with $P_{k|k-1}$ the predicted variance of $B(t_k)$. Conditional on the damage trajectory, the nuisance state is therefore integrated out exactly at linear cost, and no sampling is needed for $B$.

\subsection{Smooth Bounded Deep Discrepancy Representation}
\label{subsec:deep_discrepancy}

In high-speed precision machine tools, elastohydrodynamic lubrication (EHL) oil films thin under elevated operating temperatures, inducing localized asperity micro-slip and which can make the effective load--life exponent depart from the nominal ISO~3408-5 value $p = 3$~\cite{ISO_3408_5}. Scenario~S6 (Section~5.2) represents this structural misspecification by simulating $p = 4.5$ while the physics model assumes $p = 3$. To capture such load-exponent misspecification and unmodeled tribological dynamics without destabilizing physical monotonicity, the baseline fatigue accumulation rate is modulated by a learned multiplicative discrepancy factor $\alpha(t) = \exp\bigl(\rho(t)\bigr)$:
\begin{equation}
    \frac{\mathrm{d}D(t)}{\mathrm{d}t} = \frac{\mathrm{d}D_0(t)}{\mathrm{d}t} \cdot \exp\bigl(\rho(t)\bigr), \quad \rho(t) \in [-b, +b].
    \label{eq:total_rate}
\end{equation}
Here $\rho(t)$ represents the signed log-rate correction. In KI-SSDM, candidate discrepancy representations include:
(i) a null discrepancy $\rho \equiv 0$ (recovering the uncorrected physics backbone);
(ii) an orthogonal Basis Radial Basis Function (RBF) expansion; and
(iii) a smooth 3-layer neural network with infinitely differentiable ($C^\infty$) hyperbolic tangent activation functions and skip connections:

\begin{equation}
\label{eq:tanh_discrepancy}
\begin{aligned}
\rho(\mathbf{u},D)
={}& b\,\tanh\Bigl(
\mathbf{W}_3\,\tanh\bigl(\\
&\quad \mathbf{W}_2\,\tanh(
\mathbf{W}_1\tilde{\boldsymbol{\phi}}
+\mathbf{b}_1)\\
&\quad +\mathbf{W}_{\mathrm{skip}}
\tilde{\boldsymbol{\phi}}
+\mathbf{b}_2
\bigr)+b_3\Bigr).
\end{aligned}
\end{equation}

where $b \in \{0.30, 0.50\}$ is a declared physical authority bound (nominally $b=0.30$; $b = 0$ recovers pure physics), and $\tilde{\boldsymbol{\phi}}$ denotes normalized covariates (axial force, feed speed, instantaneous damage $D$). Because $|\tanh(\cdot)| < 1$, the multiplicative correction is strictly bounded:
\begin{equation}
    e^{-b} \le \alpha(t) \le e^b, \quad \forall t \ge 0.
    \label{eq:authority_bound}
\end{equation}
For $b = 0.30$, $\alpha(t) \in [0.741, 1.350]$. This constraint ensures strictly positive wear multipliers ($\alpha > 0$), guaranteeing that damage accumulation remains strictly monotonic ($\mathrm{d}D / \mathrm{d}t \ge 0$) while preventing explosive unphysical wear divergence.

\subsection{Linear Pool Predictive Stacking}
\label{subsec:linear_pool}

Rather than enforcing an a priori static model choice between pure physics and empirical discrepancy representations, KI-SSDM evaluates candidate representations on held-out validation units and constructs a linear pool (stacking) of predictive RUL distributions. For a given operational duty policy $\pi$, the pooled predictive density is:

\begin{equation}
\label{eq:linear_pool}
\begin{aligned}
p(\RUL \mid \mathbf{y}_{1:k})
&= (1-w)\,
p_{\mathrm{physics}}(\RUL \mid \mathbf{y}_{1:k})\\
&\quad + w\,
p_{\mathrm{discrepancy}}(
\RUL \mid \mathbf{y}_{1:k}).
\end{aligned}
\end{equation}

where predictive quantiles and intervals are computed directly from the pooled mixture CDF $F_{\mathrm{pool}}(t) = (1-w)F_{\mathrm{phys}}(t) + w F_{\mathrm{disc}}(t)$. The candidate model and stacking weight $w \in \{0, 0.25, 0.50, 0.75, 1.00\}$ are selected on held-out validation units by minimising the causal validation log-RUL loss (the mean squared error between $\log(1 + \widehat{\RUL}_{0.5}/h_0)$ and $\log(1 + \RUL/h_0)$ over causal validation forecast origins, with $h_0$ a fixed relative floor in hours); exact ties are resolved toward pure physics ($w = 0$). Across the 30 experimental fleets, model selection adapts systematically to operating conditions: under Planned duty, 8 fleets select pure physics, 11 select the Basis RBF representation, and 11 select the Deep tanh network; under Probabilistic duty, 7 fleets select pure physics, 6 select Basis RBF, and 17 select Deep tanh (12 pure and 5 stacked with physics). In well-specified scenarios (S1--S3), candidate selection distributes across representations (4 physics, 6 Basis, 5 Deep) while maintaining near-identical prognostic accuracy, demonstrating that the authority bound $b$ reliably prevents model divergence even when discrepancy capacity is activated in null regimes.

\subsection{Stochastic Process Noise: Gamma Damage Increments}
\label{subsec:gamma_noise}

To capture stochastic microstructural heterogeneity during fatigue progression, degradation increments over discrete forecasting intervals $\Delta t_k$ follow a conditional Gamma-increment process:
\begin{equation}
    \Delta D_k \sim \mathrm{Gamma}\left(\frac{\Delta t_k}{\kappa}, \; \kappa \cdot r_k\right),
    \label{eq:gamma_increments}
\end{equation}
where $\mathrm{Gamma}(\alpha, \theta)$ is parameterised in shape-scale form with shape parameter $\alpha = \Delta t_k / \kappa$ and scale parameter $\theta = \kappa \cdot r_k$. This parameterisation ensures that the conditional expectation recovers the deterministic degradation drift, $\mathbb{E}[\Delta D_k] = \alpha \theta = r_k \Delta t_k$, while the variance scales linearly with stochastic noise intensity $\kappa$, $\mathrm{Var}(\Delta D_k) = \alpha \theta^2 = \kappa r_k^2 \Delta t_k$. Here $r_k = \frac{\mathrm{d}D}{\mathrm{d}t}(t_k)$ is the instantaneous wear rate, and $\kappa \ge 0$ (in hours) governs stochastic dispersion. In the deterministic limit $\kappa \to 0$, the process recovers standard ODE numerical integration $\Delta D_k = r_k \Delta t_k$. On held-out validation units, $\kappa$ is selected from $\{0, 1, 4\}$~h under a paired one-standard-error (1-SE) safeguard: $\kappa > 0$ is adopted if and only if its improvement in the validation log-RUL loss over the null model ($\kappa = 0$) exceeds the empirical standard error of the paired score differences.

\subsection{Hierarchical Parameter Estimation and Online Edge State Filtering}
\label{subsec:hierarchical_estimation}

Because the latent damage process $D(t)$ and unit-level frailty $\eta_u$ enter the dynamics nonlinearly via Eq.~\eqref{eq:total_rate}, they cannot be marginalised in closed form like the linear-Gaussian nuisance state $B(t)$. KI-SSDM resolves this through a two-stage estimation and filtering architecture:
\begin{enumerate}[label=\textbf{Stage \arabic*:},leftmargin=3.0em]
    \item \textbf{Offline Fleet Parameter Estimation (Adaptive Quadrature Profiling)}: Fleet hyper-parameters $\boldsymbol{\Theta} = \{\mathbf{W}, \mathbf{b}, \mu_C, \sigma_C^2, \theta_B, \boldsymbol{\gamma}, \sigma_1^2, \dots, \sigma_J^2\}$ are estimated by maximizing the marginal fleet log-likelihood $\mathcal{L}(\boldsymbol{\Theta}) = \sum_{u=1}^U \log p(\mathbf{y}_u \mid \boldsymbol{\Theta})$. For each unit $u$, the unobserved random frailty $\eta_u \sim \mathcal{N}(0, \sigma_C^2)$ is integrated out numerically by Gauss--Hermite quadrature with $Q = 25$ nodes~\cite{Liu_1994, Pinheiro_1995}:
    
\begin{equation}
\label{eq:quadrature_marginal}
\begin{aligned}
p(\mathbf{y}_u \mid \boldsymbol{\Theta})
&= \int_{-\infty}^{\infty}
p\bigl(\mathbf{y}_u \mid
\mathbf{D}_u(\eta_u),
\hat{\boldsymbol{\beta}},
\boldsymbol{\Theta}\bigr)\\
&\quad \times
\frac{\exp\!\left(
-\eta_u^2/(2\sigma_C^2)
\right)}
{\sqrt{2\pi}\sigma_C}
\,\mathrm{d}\eta_u\\[2pt]
&\approx \frac{1}{\sqrt{\pi}}
\sum_{q=1}^{Q} w_q\\
&\quad \times
p\bigl(\mathbf{y}_u \mid
\mathbf{D}_u(\sqrt{2}\sigma_C z_q),\\
&\qquad
\hat{\boldsymbol{\beta}},
\boldsymbol{\Theta}\bigr).
\end{aligned}
\end{equation}

    where $z_q$ and $w_q$ denote the standard Gauss-Hermite quadrature nodes and weights ($Q = 25$). At each quadrature node, the deterministic damage trajectory $\mathbf{D}_u(\eta_u)$ is evaluated via RK4 integration of Eq.~\eqref{eq:total_rate}. Conditioning on $\mathbf{D}_u(\eta_u)$, emission weights $\boldsymbol{\beta}$ are profiled out in closed form across all units via the pooled GLS estimator in Eq.~\eqref{eq:gls_solution}. The conditional density $p(\mathbf{y}_u \mid \mathbf{D}_u(\eta_u), \hat{\boldsymbol{\beta}}, \boldsymbol{\Theta})$ is then evaluated in exact $O(K_u)$ time via the Kalman innovations representation in Eq.~\eqref{eq:kalman_innovations}. Fleet hyper-parameters $\boldsymbol{\Theta}$ are optimised by a projected quasi-Newton (BFGS) method with central finite-difference gradients, an Armijo line search, and box and $\ell_1$-ball constraints, alternated with closed-form updates of the emission and noise parameters, until strict stationarity criteria are satisfied ($\|\nabla \mathcal{L}\| \le 5 \times 10^{-5}$ for Exitflag~1, or $< 5 \times 10^{-4}$ for Exitflag~2).
    \item \textbf{Online Filtering and RUL Prognosis}: In online deployment, all fleet parameters $\hat{\boldsymbol{\Theta}}$ and emission weights $\hat{\boldsymbol{\beta}}$ remain frozen. At each observation window, the damage state advances by Eq.~\eqref{eq:total_rate}, the frailty posterior is carried on an adaptive support of weighted points, and, conditional on each support point, the nuisance state is updated by a Kalman step on the innovation $\mathbf{y}_k - \hat{\mathbf{y}}_{k|k-1}$. At each forecast origin, the RUL distribution is obtained by simulating 384 forward damage paths in each of four independent replicates under the planned or probabilistic duty model to the threshold $\Dcrit = 1$, pooling the physics and discrepancy predictive distributions with weight $w$, and applying the conformal interval of Section~\ref{sec:conformal_uq}.
\end{enumerate}

Algorithm~\ref{alg:hierarchical_kissdm} summarises the end-to-end mathematical execution sequence.

\begin{algorithm}[!t]
\caption{Two-Stage Hierarchical Estimation and Edge State Filtering in KI-SSDM.}
\label{alg:hierarchical_kissdm}
\begin{algorithmic}[1]
\Require Training fleet trajectories $\{\mathbf{y}_u, \mathbf{u}_u, \bPhi_u\}_{u=1}^U$; test stream $\bPhi_k \in \R^{23}$; miscoverage level $\alpha = 0.10$.
\Ensure Point forecast $\widehat{\RUL}_{0.5}$, conformal predictive interval $\mathcal{I}_{1-\alpha}(t_k)$.
\Statex \textbf{Stage 1: Offline Fleet Profiling \& optimisation}
\Repeat
    \For{each training unit $u \in \{1, \dots, U\}$}
        \For{each Gauss-Hermite quadrature node $q \in \{1, \dots, Q\}$}
            \State Set frailty $\eta_{u,q} = \sqrt{2}\sigma_C z_q$;
            \State Integrate $\mathbf{D}_{u,q}$ via RK4 using Eq.~\eqref{eq:total_rate};
        \EndFor
    \EndFor
    \State Evaluate pooled emission weights $\hat{\boldsymbol{\beta}}$ analytically via GLS (Eq.~\eqref{eq:gls_solution});
    \For{each unit $u$}
        \State Compute $p(\mathbf{y}_u \mid \mathbf{D}_{u,q}, \hat{\boldsymbol{\beta}})$ via Kalman innovations (Eq.~\eqref{eq:kalman_innovations});
        \State Accumulate quadrature marginal likelihood $p(\mathbf{y}_u \mid \boldsymbol{\Theta})$ via Eq.~\eqref{eq:quadrature_marginal};
    \EndFor
    \State Update $\boldsymbol{\Theta}$ by a projected BFGS step on $-\sum_u \log p(\mathbf{y}_u \mid \boldsymbol{\Theta})$ (finite-difference gradients, Armijo line search);
\Until{$\|\nabla \mathcal{L}\| < 5 \times 10^{-5}$ (Exitflag 1) or $< 5 \times 10^{-4}$ (Exitflag 2)};
\State optimise stacking weight $w \in \{0, 0.25, 0.50, 0.75, 1.00\}$ on validation units via the log-RUL loss (exact ties to $w=0$);
\State Compute origin-stratified calibration quantiles $\hat{q}_{1-\alpha}(t_0) = s_{(k)}$ with $k = \lceil (n+1)(1-\alpha) \rceil$;
\Statex \textbf{Stage 2: Online Filtering and RUL Prognosis}
\State Ingest $\bPhi_k$; extract torque innovation $e_k = \tau_k - h_\psi(\mathbf{r}_k)$;
\State Compute $\rho_k = \rho(\mathbf{u}_k, D_{k-1})$ via $C^\infty$ tanh network (Eq.~\eqref{eq:tanh_discrepancy});
\State Advance damage state: $D_k = D_{k-1} + \frac{\mathrm{d}D_0}{\mathrm{d}t} e^{\rho_k} \Delta t_k$;
\State Recursive Kalman update on 1D nuisance state $\hat{B}_{k|k}$ using innovation $\mathbf{y}_k - \mathbf{X}_k(D_k) \hat{\boldsymbol{\beta}}$;
\State Project RUL CDF $F_{\mathrm{pool}}(t)$ by evaluating prospective threshold crossings $D(t) \ge \Dcrit = 1.0$;
\State Extract point forecast $\widehat{\RUL}_{0.5} = F_{\mathrm{pool}}^{-1}(0.5)$;
\State Construct conformal interval $\mathcal{I}_{1-\alpha}(t_k) = \bigl[\max(0, \widehat{\RUL}_{0.5} - \hat{q}_{1-\alpha}), \; \widehat{\RUL}_{0.5} + \hat{q}_{1-\alpha}\bigr]$.
\end{algorithmic}
\end{algorithm}

% ========================  4. CONFORMAL UQ  ========================
\section{Origin-Stratified Conformal Prediction}
\label{sec:conformal_uq}

Predictive maintenance in industrial manufacturing requires reliable uncertainty bounds with provable statistical coverage. Traditional Bayesian credible intervals rely heavily on correct prior and likelihood specifications, often under-covering true failure times when operational regimes shift. Conformal prediction provides distribution-free, finite-sample predictive guarantees under the sole mathematical assumption of unit exchangeability~\cite{Angelopoulos_2023, Zhang_CP_2025}.

KI-SSDM implements origin-stratified split-conformal prediction. Let $\mathcal{D}_{\mathrm{cal}}(t_0) = \{u_1, \dots, u_n\}$ denote the set of $n$ surviving, eligible calibration units at a prespecified forecast origin $t_0$. The nonconformity score for unit $u \in \mathcal{D}_{\mathrm{cal}}(t_0)$ is defined by the absolute median RUL prediction error:
\begin{equation}
    s_u(t_0) = \bigl|\RUL_u(t_0) - \widehat{\RUL}_{u,0.5}(t_0)\bigr|,
    \label{eq:nonconformity_score}
\end{equation}
where $\widehat{\RUL}_{u,0.5}(t_0)$ is the median RUL forecast generated by the pooled predictive CDF.

Given a declared miscoverage rate $\alpha \in (0, 1)$ (nominally $\alpha = 0.10$ for 90\% coverage), let $s_{(1)} \le s_{(2)} \le \dots \le s_{(n)}$ denote the sorted calibration nonconformity scores. The origin-stratified conformal quantile is evaluated via the exact finite-sample order statistic:
\begin{equation}
    k = \lceil (n + 1)(1 - \alpha) \rceil, \quad \hat{q}_{1-\alpha}(t_0) = \begin{cases} s_{(k)}, & \text{if } k \le n, \\ \infty, & \text{if } k > n. \end{cases}
    \label{eq:conformal_quantile}
\end{equation}
For $n = 24$ calibration units per fleet at $\alpha = 0.10$, $k = \lceil 25 \times 0.90 \rceil = 23$, giving an exact finite-sample continuous-score rank coverage of $k / (n + 1) = 23 / 25 = 92\%$.

For any unseen test unit $u_{\mathrm{test}}$ evaluated at origin $t_0$, the calibrated predictive interval is formed by intersecting with $[0, \infty)$ to enforce physical non-negativity:

\begin{equation}
\label{eq:conformal_interval}
\mathcal{I}_{1-\alpha}(t_0)
=\bigl[L_\alpha(t_0),U_\alpha(t_0)\bigr],
\end{equation}
\begin{equation}
\begin{aligned}
L_\alpha(t_0)
&=\max\bigl\{0,\,
\widehat{\RUL}_{0.5}(t_0)\\
&\qquad-\hat{q}_{1-\alpha}(t_0)\bigr\},\\
U_\alpha(t_0)
&=\widehat{\RUL}_{0.5}(t_0)
+\hat{q}_{1-\alpha}(t_0).
\end{aligned}
\end{equation}


By exchangeability of units across calibration and test allocations at origin $t_0$, the marginal coverage satisfies the finite-sample guarantee:
\begin{equation}
    \mathbb{P}\Bigl(\RUL_{\mathrm{test}}(t_0) \in \mathcal{I}_{1-\alpha}(t_0)\Bigr) \ge 1 - \alpha.
    \label{eq:coverage_guarantee}
\end{equation}

When evaluating interval prediction quality, empirical coverage alone is insufficient: an uninformative, arbitrarily wide interval trivially achieves 100\% empirical coverage. A rigorous comparative assessment requires evaluating \emph{sharpness} (mean prediction interval width $\overline{\Delta}_{\mathcal{I}} = U - L$) at equal nominal coverage ($1 - \alpha = 0.90$), alongside the Winkler Interval Score $\mathrm{IS}_\alpha$~\cite{Winkler_1972, Gneiting_2007}:
\begin{equation}
    \mathrm{IS}_\alpha(\mathcal{I}, y) = (U - L) + \frac{2}{\alpha}(L - y)\mathbf{1}[y < L] + \frac{2}{\alpha}(y - U)\mathbf{1}[y > U],
    \label{eq:interval_score}
\end{equation}
where $\mathcal{I} = [L, U]$ and $y$ is the ground-truth failure time. The interval score rewards sharpness while penalizing miscoverage outliers with slope $2/\alpha = 20$. Origin stratification calibrates the nonconformity score independently at each inspection origin, allowing the interval width to contract dynamically as the component approaches retirement.

% ========================  5. EXPERIMENTAL SETUP  ========================
\section{Multi-Fleet Experimental Benchmark Setup}
\label{sec:experiments}

\subsection{Multi-Fleet Benchmark Hierarchy}
\label{subsec:fleet_hierarchy}

To evaluate KI-SSDM under realistic, non-cherry-picked conditions, a multi-fleet synthetic benchmark was constructed (Table~\ref{tab:benchmark_setup}). The simulator uses the fatigue law of Eq.~\eqref{eq:miner_damage} with log-normal frailty (wear rate proportional to $F^p v\,e^{\beta_T T + \eta_u}$) in dimensionless units: it is not calibrated to the physical screw of Section~\ref{sec:bench_features}, and simulated time in hours is a nominal benchmark scale rather than an equivalent service life. Each unit is observed every 8~h on $J = 3$ abstract telemetry channels whose means are linear in $D$, $B$ and the log-load, log-speed and temperature covariates (channel~1 also contains a $D^2$ term), as in Eq.~\eqref{eq:observation_model}; the 500~ms windows, the digital twin and the 23-dimensional feature vector $\bPhi_k$ of Section~\ref{sec:bench_features} are not used in the benchmark. Across the 480 held-out test units, the simulated mean component lifetime is $371.9$~h (unit standard deviation $78.2$~h; fleet means spanning $328.7\text{--}443.3$~h, fleet standard deviation $36.1$~h). In this context, an MAE of $16.40$~h under planned duty represents only $4.4\%$ of component lifetime, and $18.29$~h under probabilistic duty represents $4.9\%$ of component lifetime.

The benchmark simulates $N_{\mathrm{fleets}} = 30$ independent machine fleets ($N_{\mathrm{units}} = 2{,}220$ total simulated units, partitioned into 24 training, 10 validation, 24 calibration, and 16 held-out test units per fleet, yielding $16 \times 30 = 480$ test units overall) across two operational duty protocols:
\begin{itemize}[leftmargin=1.5em]
    \item \textbf{Planned Duty Protocol}: The future operational profile (speed, axial load, reversal frequency) is known deterministically over the entire remaining life horizon.
    \item \textbf{Probabilistic Duty Protocol}: Future duty cycles are non-stationary Markovian sequences sampled from operational duty transition matrices, reflecting real-world schedule uncertainty.
\end{itemize}

\begin{table}[pos=!t]
\centering
\caption{Hierarchical benchmark experimental configuration.}
\label{tab:benchmark_setup}
\renewcommand{\arraystretch}{1.12}
\setlength{\tabcolsep}{3.5pt}
\small
\begin{tabularx}{\columnwidth}{@{}l X@{}}
\toprule
\textbf{Configuration Parameter} & \textbf{Value / Specification} \\
\midrule
Total Fleets / Total Units & 30 independent fleets / 2,220 units \\
Unit Allocation per Fleet  & 24 Train / 10 Validation / 24 Cal / 16 Test \\
Operating Policies         & Planned Duty, Probabilistic Markov Duty \\
Units and Scale & Dimensionless; not calibrated to the physical screw \\
Mean Test-Unit Life & $371.9 \pm 78.2$~h (480 test units) \\
Sampling Interval & 8~h; $J = 3$ abstract channels \\
Failure Threshold ($\Dcrit$)  & $D = 1.0$ (Spall initiation / raceway pitting) \\
Measurement Noise ($\sigma_j$) & 0.045 per channel (0.16 in S5) \\
OU Mean Reversion ($\theta_B$) & Simulated $1/18 \approx 0.056~\mathrm{h}^{-1}$ ($t_{1/2} = 12.5$~h); estimated in $\boldsymbol{\Theta}$ \\
Conformal Nominal Level ($1-\alpha$) & 90\% marginal coverage \\
Optimisation Stationarity  & Exitflag~1: $\|\nabla\| < 5 \times 10^{-5}$; Exitflag~2: $< 5 \times 10^{-4}$ \\
\bottomrule
\end{tabularx}
\end{table}

Each fleet is strictly partitioned into disjoint subsets: 24 training units, 10 validation units, 24 calibration units, and 16 test units. Information leakage is strictly prevented: calibration and test units are completely inaccessible during parameter estimation and model selection.

\subsection{Degradation Scenarios}
\label{subsec:scenarios}

The 30-fleet evaluation spans six distinct operational degradation regimes ($5$ independent fleet replicates per scenario, totaling $30$ fleets):
\begin{enumerate}[label=\textbf{S\arabic*.},leftmargin=2.2em]
    \item \textbf{Matched Null Case} (\texttt{fatigue\_matched}): The true simulator exactly follows baseline physical fatigue ($p = 3.0$, zero unmodeled wear). Tests whether the model overfits or corrupts a known physics model.
    \item \textbf{Accelerated Discrepancy} (\texttt{fatigue\_discrepancy}): The true wear rate is multiplied by an unmodelled smooth factor $\exp\bigl(0.28\sin(3D) + 0.10\,D\tanh T\bigr)$ over the whole life.
    \item \textbf{Emission Distortion} (\texttt{fatigue\_emission\_mismatch}): Sensor telemetry emission mapping exhibits unmodelled cubic distortion ($+0.15D^3$ on channel~1, $-0.08D^3$ on channel~3), testing robustness against sensor calibration drift.
    \item \textbf{Combined Mismatch} (\texttt{fatigue\_combined\_mismatch}): The S2 rate discrepancy and the S3 emission distortion together.
    \item \textbf{Telemetry Dropouts} (\texttt{fatigue\_weak\_gap}): Telemetry is blanked between 120 and 240~h, sensor noise is raised to 0.16, and the S2 rate discrepancy is present, testing filtering robustness under missing and noisy data.
    \item \textbf{Structural Load Misspecification} (\texttt{fatigue\_load\_exponent}): The true mechanical load exponent is $p = 4.5$, whereas the nominal physics model assumes $p = 3.0$. Creates severe, duty-dependent log-rate discrepancies.
\end{enumerate}

\subsection{Baseline Comparators}
\label{subsec:baselines}

KI-SSDM is benchmarked against five established methods in industrial prognostics:
\begin{itemize}[leftmargin=1.5em]
    \item \textbf{Physics SSM}: Pure deterministic state-space model following baseline fatigue laws without learned corrections ($w = 0$).
    \item \textbf{Basis Hybrid}: Physics backbone with a bounded log-rate discrepancy built from 5 Gaussian RBF centres in $D$ plus 5 linear covariate terms, fitted by the same marginal likelihood.
    \item \textbf{Bayesian Particle Filter (PF)}: 1,200 particles with systematic resampling, sequentially updating the fatigue state with empirical observation noise covariance.
    \item \textbf{Inverse PINN}: A small tanh network (3 and 2 hidden units) in the fatigue rate with trainable wear coefficients and physics-residual weight 1, fitted by the same marginal-likelihood optimiser and iteration budget as the deep discrepancy.
    \item \textbf{Deep LSTM}: Purely data-driven single-layer LSTM (hidden size 16 or 32, selected on validation) followed by a ReLU fully connected head, trained with Adam (learning rate 0.001 or 0.003, at most 100 epochs, batch 16, early stopping with patience 8), mapping telemetry window sequences directly to RUL.
\end{itemize}

\subsection{Cross-Domain Benchmark on Run-to-Failure Bearings (PRONOSTIA)}
\label{subsec:pronostia_benchmark_setup}

To evaluate algorithmic generalizability on external vibration benchmarks where run-to-failure testing is standard, the proposed KI-SSDM methodology was independently evaluated on the public IEEE PHM 2012 / PRONOSTIA bearing platform~\cite{Nectoux_2012} (17 run-to-failure bearings across three operating regimes). On this benchmark KI-SSDM was significantly more accurate than exponential Bayesian and particle-filter baselines (Holm $p \le 0.0013$, 17/17 bearings) but did not differ from a deep LSTM ($p = 0.50$) or from a signal-free constant-RUL predictor ($p = 0.67$); the vibration features available do not support accurate bearing prognosis by any of the methods tested (Supplementary Annexure).
% ========================  6. RESULTS  ========================
\section{Experimental Results and Detailed Analysis}
\label{sec:results}

\subsection{Overall Benchmark Prognostic Accuracy}
\label{subsec:overall_results}

Table~\ref{tab:overall_results} presents the comparative performance across all 30 independent fleets ($N = 30$, spanning 6 distinct operational scenarios $\times$ 5 fleet replicates) under both Planned and Probabilistic duty policies. The prognostic gain is defined as:

\begin{equation}
\begin{aligned}
\mathrm{Gain} ={}&
\MAE_{\mathrm{Comparator}} \\
&- \MAE_{\mathrm{KI\text{-}SSDM}}.
\end{aligned}
\label{eq:gain_def}
\end{equation}
\noindent
Positive gain values indicate superior performance
of the proposed KI-SSDM.


%========================================================
% USER-CONFIGURABLE COLUMN WIDTHS FOR TABLE 5 (tab:overall_results)
% Adjust each column independently.
% Use distinct names to avoid conflicts with other tables.
%========================================================
\newlength{\ovColA}
\newlength{\ovColB}
\newlength{\ovColC}
\newlength{\ovColD}
\newlength{\ovColE}
\newlength{\ovColF}
\newlength{\ovColG}
\newlength{\ovColH}
\newlength{\ovColI}

\setlength{\ovColA}{1.85cm} % Policy
\setlength{\ovColB}{2.30cm} % Method / Comparator
\setlength{\ovColC}{1.05cm} % MAE (h)
\setlength{\ovColD}{1.25cm} % Gain vs KI
\setlength{\ovColE}{2.25cm} % 95% CI [h]
\setlength{\ovColF}{1.85cm} % Win/Loss/Tie
\setlength{\ovColG}{1.85cm} % pt / pW
\setlength{\ovColH}{1.15cm} % pHolm
\setlength{\ovColI}{2.50cm} % Statistical Verdict

\begin{table*}[pos=!t]
\centering
\caption{Comprehensive prognostic benchmark performance across 30 independent fleets ($N=30$). Proposed KI-SSDM serves as the reference benchmark ($\mathrm{MAE} = 16.40$~h Planned; $18.29$~h Probabilistic). Gains ($\Delta\mathrm{MAE} = \mathrm{MAE}_{\mathrm{Comparator}} - \mathrm{MAE}_{\mathrm{KI\text{-}SSDM}}$) and 95\% paired confidence intervals are reported in hours; two-sided paired $t$-test $p$-values ($p_t$), exact Wilcoxon signed-rank $p$-values ($p_W$), and Holm-adjusted $p$-values ($p_{\mathrm{Holm}}$, five comparisons per duty~\cite{Holm_1979}) are computed across independent fleets. Ties are fleet-level $|\Delta\mathrm{MAE}| \le 0.05$~h.}
\label{tab:overall_results}
\renewcommand{\arraystretch}{1.15}
\setlength{\tabcolsep}{2.5pt}
\small
\begin{tabular}{@{}
>{\raggedright\arraybackslash}p{\ovColA}
>{\raggedright\arraybackslash}p{\ovColB}
>{\centering\arraybackslash}p{\ovColC}
>{\centering\arraybackslash}p{\ovColD}
>{\centering\arraybackslash}p{\ovColE}
>{\centering\arraybackslash}p{\ovColF}
>{\centering\arraybackslash}p{\ovColG}
>{\centering\arraybackslash}p{\ovColH}
>{\raggedright\arraybackslash}p{\ovColI}
@{}}
\toprule
\textbf{Policy} & \textbf{Method / Comparator} & \textbf{MAE (h)} & \textbf{Gain vs.\ KI} & \textbf{95\% CI [h]} & \textbf{Win/Loss/Tie} & $\bm{p_t}$ \textbf{/} $\bm{p_W}$ & $\bm{p_{\mathrm{Holm}}}$ & \textbf{Statistical Verdict} \\
\midrule
\multirow{6}{=}{\textbf{Planned Duty}}
& \textbf{KI-SSDM (Proposed)} & \textbf{16.40} & \textbf{--- (Ref)} & \textbf{---} & \textbf{---} & \textbf{---} & \textbf{---} & \textbf{Reference Benchmark} \\
& Inverse PINN                & 17.17          & \textbf{+0.77}     & $[+0.31, \; +1.24]$  & \textbf{22 / 8 / 0}      & \textbf{0.0020 / 0.0012} & \textbf{0.0079} & \textbf{Significant ($p_{\mathrm{adj}} < 0.05$)} \\
& Bayesian PF                 & 17.38          & \textbf{+0.99}     & $[+0.04, \; +1.93]$  & \textbf{17 / 8 / 5}      & 0.0411 / 0.0293          & 0.1073          & n.s.\ after Holm \\
& Physics SSM                 & 17.41          & \textbf{+1.02}     & $[+0.07, \; +1.96]$  & \textbf{16 / 6 / 8}      & 0.0358 / 0.0138          & 0.1073          & n.s.\ after Holm \\
& Basis Hybrid                & 16.65          & +0.26              & $[-0.38, \; +0.90]$  & 8 / 8 / 14               & 0.4175 / 0.7337          & 0.4175          & Not significant \\
& Deep LSTM                   & 29.41          & \textbf{+13.01}    & $[+10.42, \; +15.61]$& \textbf{30 / 0 / 0}      & $\bm{<10^{-10}}$ \textbf{/} $\bm{<10^{-8}}$ & $\bm{<10^{-9}}$  & \textbf{Significant ($p_{\mathrm{adj}} < 0.05$)} \\
\midrule
\multirow{6}{=}{\textbf{Probabilistic Duty}}
& \textbf{KI-SSDM (Proposed)} & \textbf{18.29} & \textbf{--- (Ref)} & \textbf{---} & \textbf{---} & \textbf{---} & \textbf{---} & \textbf{Reference Benchmark} \\
& Inverse PINN                & 18.81          & \textbf{+0.52}     & $[+0.04, \; +1.00]$  & \textbf{18 / 10 / 2}     & 0.0341 / 0.0345          & 0.1364          & n.s.\ after Holm \\
& Bayesian PF                 & 18.85          & +0.56              & $[-0.13, \; +1.25]$  & 16 / 9 / 5               & 0.1099 / 0.1191          & 0.2904          & Not significant \\
& Physics SSM                 & 18.87          & +0.58              & $[-0.11, \; +1.27]$  & 15 / 6 / 9               & 0.0968 / 0.1074          & 0.2904          & Not significant \\
& Basis Hybrid                & 18.47          & +0.18              & $[-0.18, \; +0.54]$  & 10 / 6 / 14              & 0.3163 / 0.3300          & 0.3163          & Not significant \\
& Deep LSTM                   & 29.82          & \textbf{+11.53}    & $[+9.41, \; +13.66]$ & \textbf{30 / 0 / 0}      & $\bm{<10^{-11}}$ \textbf{/} $\bm{<10^{-8}}$ & $\bm{<10^{-10}}$ & \textbf{Significant ($p_{\mathrm{adj}} < 0.05$)} \\
\bottomrule
\end{tabular}
\end{table*}
Under the Planned Duty policy, KI-SSDM achieves an overall fleet-mean MAE of \textbf{16.40~h}, outperforming the deep LSTM network (29.41~h, representing a +13.01~h gain, $p_t < 10^{-10}$, $p_{\mathrm{Holm}} < 10^{-9}$, winning 30 out of 30 fleets). Against hybrid models, KI-SSDM achieves statistically significant prognostic gains over Inverse PINN (+0.77~h, 95\% CI $[+0.31, +1.24]$~h, $p_t = 0.0020$, $p_{\mathrm{Holm}} = 0.0079$, winning 22 of 30 fleets). Against uncorrected physical models, KI-SSDM achieves unadjusted paired $t$-test gains over Physics SSM (+1.02~h, raw $p_t = 0.0358$, Wilcoxon $p = 0.0138$) and Bayesian PF (+0.99~h, raw $p_t = 0.0411$, Wilcoxon $p = 0.0293$). After Holm correction across the five comparators within each duty, these differences are not significant ($p_{\mathrm{Holm}} = 0.107$). Outside S6 the bounded correction changes forecasts little, whichever candidate validation selects, so KI-SSDM stays close to the physics baselines (mean non-S6 gain $+0.02$~h), while the dominant aggregate gain is concentrated under structural load-exponent misspecification in Scenario~S6 (+5.99~h, within-scenario Holm $p = 0.0277$, accounting for 98.2\% of the total +6.10~h planned physics gain). Under the highly volatile Probabilistic Duty regime, KI-SSDM achieves a mean MAE of \textbf{18.29~h}, is more accurate than the LSTM baseline by +11.53~h ($p_{\mathrm{Holm}} < 10^{-10}$), and does not differ significantly from the physics-based and hybrid baselines.

\subsection{Continuous Latent State-Space Tracking}
\label{subsec:state_tracking}

Fig.~\ref{fig:states_filter} illustrates joint latent state tracking and operational residual decoupling on a representative degraded benchmark unit across a six-panel diagnostic grid.

In the top row, panel~(a) shows the latent damage $D(t)$ accumulating monotonically, with the smoothed and filtered estimates following the simulated truth. Panel~(b) shows the standardised Ornstein-Uhlenbeck nuisance state $B(t)$ (simulated with $\theta_B = 1/18\,\si{h^{-1}}$) and its posterior band; its excursions are reversible and are not absorbed into the damage path. Panel~(c) shows channel-1 telemetry with the fitted emission mean and its 90\% measurement-noise band ($\pm 1.645\,\sigma_1$).

In the bottom row, panel~(d) shows how the effective number of points in the online adaptive frailty grid shrinks as evidence accumulates over the unit life (from 73.8 to 4.6; the 25 Gauss--Hermite nodes are used only in offline training), and panel~(e) shows the frailty support points $\eta_u$ at the final origin. Panel~(f) shows the standardised pre-update innovations of the three telemetry channels, which fluctuate around zero without visible trend.

\begin{figure*}[pos=!t]
  \centering
  \includegraphics[width=0.98\textwidth, keepaspectratio]{E01_States_filter.png}
  \caption{Latent-state filtering for test unit fatigue\_matched\_021 (Experiment~1, S1). (a) Damage $D(t)$: simulated truth (black), smoothed and filtered estimates. (b) Standardised Ornstein--Uhlenbeck nuisance state $B(t)$ (simulated with $\theta_B = 1/18\,\mathrm{h}^{-1}$) with its posterior band. (c) Channel-1 telemetry (dots) and fitted emission mean with 90\% measurement-noise band ($\pm1.645\,\sigma_1$). (d) Effective number of points in the online adaptive frailty grid along the unit life. (e) Frailty support points $\eta_u$ at the final origin. (f) Standardised pre-update innovations for the three telemetry channels.}
  \label{fig:states_filter}
\end{figure*}

\subsection{Optimisation Stationarity and Discrepancy Convergence}
\label{subsec:optim_convergence}

Numerical convergence of the multi-start deep discrepancy encoder progresses through primary encoder identification followed by output-coordinate stationarity iterations under a two-tier convergence rule:
\begin{itemize}[leftmargin=1.5em]
    \item \textbf{Exitflag 1 (KKT Stationarity)}: Achieved when the projected gradient norm falls below $5 \times 10^{-5}$.
    \item \textbf{Exitflag 2 (Numerical Resolution Limit)}: Triggered when the projected gradient norm is below $5 \times 10^{-4}$ and either no further descent is possible after a $4\times$ smaller gradient step refresh, or the objective stalls within a relative $10^{-8}$ window over 10 iterations.
\end{itemize}
This dual classification ensures that models converging within practical machine float precision are declared eligible for validation scoring rather than being erroneously aborted on iteration budgets. Across the multi-fleet benchmark evaluations, candidate models reliably satisfied these first-order stationarity criteria before entering the validation-gated linear pooling stage.

\subsection{RUL Prognosis under Planned and Probabilistic Duty}
\label{subsec:prognosis_bands}

Fig.~\ref{fig:prognosis_combined} shows RUL forecasts for one S6 fleet (Experiment~6, load-exponent misspecification), in which validation selected the stacked deep discrepancy under both duties. Under planned duty KI-SSDM reduced the fleet MAE from 24.24~h (Physics SSM and Basis Hybrid) to 15.36~h, with 88.7\% calibrated coverage at a mean band width of 81.6~h against 93.5\% at 133.6~h for Physics SSM. Inverse PINN was more accurate in this fleet (14.39~h). Under probabilistic duty the corresponding MAEs are 18.66~h (KI-SSDM), 23.08~h (Physics SSM) and 15.88~h (Inverse PINN), with KI-SSDM coverage 90.3\% at 104.7~h. In three of the five S6 fleets the basis correction, not the deep one, was selected.

\begin{figure*}[pos=!t]
  \centering
  \includegraphics[width=0.98\textwidth, keepaspectratio]{Fig6_Prognosis_Combined.png}
  \caption{RUL prognosis for one S6 fleet (Experiment~6: load-exponent misspecification, replicate~1), in which validation selected the stacked deep discrepancy ($w = 0.75$ planned, $0.50$ probabilistic). Panels per row: (i) forecasts for test unit fatigue\_load\_exponent\_021 at calendar origins 80, 160, 240 and 320~h with the KI-SSDM 90\% conformal band (widths 151, 130, 66 and 73~h planned; 155, 124, 78 and 83~h probabilistic), which misses the true RUL at 240~h under both duties; (ii) predicted against true RUL for the fleet's 62 unit-origin forecasts (16 test units); (iii) share of forecasts within $\pm$20\% of true RUL by life quarter; (iv) row coverage of the raw, fleet-calibrated and unit-calibrated bands; (v) mean calibrated band width; (vi) forecasts with the discrepancy on and off under identical posterior draws. (a) Planned duty: fleet MAE 15.36~h (KI-SSDM), 24.24~h (Physics SSM, Basis Hybrid), 14.39~h (Inverse PINN); KI-SSDM calibrated coverage 88.7\% at mean width 81.6~h. (b) Probabilistic duty: MAE 18.66~h, 23.08~h and 15.88~h; coverage 90.3\% at 104.7~h. Physics SSM, Basis Hybrid and Bayesian PF forecasts nearly coincide and overlap in the plots.}
  \label{fig:prognosis_combined}
\end{figure*}

\subsection{Comparative Benchmark MAE Across All Scenarios}
\label{subsec:benchmark_mae_analysis}

Fig.~\ref{fig:benchmark_mae} visualizes the multi-fleet MAE distribution across all methods. The LSTM shows both higher errors and a wider spread across fleets; the model-based methods, KI-SSDM included, have similar spreads within each scenario.

\begin{figure*}[pos=!t]
  \centering
  \includegraphics[width=0.96\textwidth, keepaspectratio]{Fig1_MAE_by_scenario.png}
  \caption{Fleet-level RUL prognostic accuracy benchmark across all 30 independent simulated fleets (Planned Duty, left; Probabilistic Duty, right). Boxes indicate median and interquartile ranges; dashed whiskers show full fleet spreads; red crosses denote outlier fleets.}
  \label{fig:benchmark_mae}
\end{figure*}

\subsection{Hierarchical Population Inference and Parameter Identifiability}
\label{subsec:identifiability_audit}

Fig.~\ref{fig:hierarchical_ident} presents the four-panel hierarchical fleet frailty and latent state identifiability audit across training units.


In the top row, panel~(a) presents the posterior unit-specific frailty estimates
($\hat{\eta}_u \pm 2\,\mathrm{SD}$) for the 24 training units, showing unit-to-unit heterogeneity around the population prior $\mathcal{N}(0,\sigma_C^2)$. Panel~(b) shows the relative damage uncertainty, $\mathrm{SD}(D)/\max\{D,0.01\}$, against the estimated damage $D$; it decreases as damage accumulates and stays below the 0.25 early-warning level.


In the bottom row, panel~(c) shows the posterior frailty standard deviation $\sigma(\eta \mid y)$ against estimated damage, falling from 0.149 to 0.023 as operational evidence accumulates, with the 0.10 early-warning level marked. Panel~(d) shows the paired RUL change ($\pm$ Monte-Carlo SD, hours) for the five declared audit contrasts; one contrast is slightly negative ($-0.16\,\si{h}$).

\begin{figure*}[pos=!t]
  \centering
  \includegraphics[width=0.92\textwidth, keepaspectratio]{E01_Hierarchical_inference.png}
  \caption{Frailty and damage identifiability for Fleet~1 (Experiment~1, S1). (a) Posterior frailty mean $\pm 2$\,SD for the 24 training units. (b) Relative damage uncertainty $\mathrm{SD}(D)/\max(D,0.01)$ against estimated damage, with the 0.25 early-warning level. (c) Posterior frailty SD against estimated damage (0.149 to 0.023), with the 0.10 early-warning level. (d) Paired RUL change ($\pm$ Monte-Carlo SD) for the five declared audit contrasts.}
  \label{fig:hierarchical_ident}
\end{figure*}

\subsection{Scenario Breakdown and Physical Misspecification}
\label{subsec:scenario_breakdown}

Table~\ref{tab:scenario_breakdown} presents the scenario-specific MAE breakdown across the six evaluation regimes.

%========================================================
% USER-CONFIGURABLE COLUMN WIDTHS FOR TABLE 6 (tab:scenario_breakdown)
% Adjust each column independently.
% Use distinct names to avoid conflicts with other tables.
%========================================================
\newlength{\scColA}
\newlength{\scColB}
\newlength{\scColC}
\newlength{\scColD}
\newlength{\scColE}
\newlength{\scColF}
\newlength{\scColG}
\newlength{\scColH}

\setlength{\scColA}{3.00cm} % Scenario
\setlength{\scColB}{1.65cm} % Duty Policy
\setlength{\scColC}{1.35cm} % Physics SSM
\setlength{\scColD}{1.35cm} % Bayesian PF
\setlength{\scColE}{1.35cm} % Inverse PINN
\setlength{\scColF}{0.95cm} % LSTM
\setlength{\scColG}{1.35cm} % KI-SSDM
\setlength{\scColH}{2.70cm} % KI Gain vs Physics
\newlength{\scColI}
\setlength{\scColI}{1.35cm} % Basis Hybrid

\begin{table*}[pos=!t]
\centering
\caption{Comprehensive scenario-specific MAE breakdown (hours) across six evaluation regimes (five independent fleet replicates per scenario, totaling 30 fleets). The proposed KI-SSDM demonstrates competitive performance across the evaluated regimes, with particularly strong improvements under load-exponent misspecification (S6).}
\label{tab:scenario_breakdown}
\renewcommand{\arraystretch}{1.15}
\setlength{\tabcolsep}{3pt}
\small
\begin{tabular}{@{}
>{\raggedright\arraybackslash}p{\scColA}
>{\raggedright\arraybackslash}p{\scColB}
>{\centering\arraybackslash}p{\scColC}
>{\centering\arraybackslash}p{\scColD}
>{\centering\arraybackslash}p{\scColE}
>{\centering\arraybackslash}p{\scColI}
>{\centering\arraybackslash}p{\scColF}
>{\centering\arraybackslash}p{\scColG}
>{\raggedright\arraybackslash}p{\scColH}
@{}}
\toprule
\textbf{Scenario}
& \textbf{Duty Policy}
& \textbf{Physics SSM}
& \textbf{Bayesian PF}
& \textbf{Inverse PINN}
& \textbf{Basis Hybrid}
& \textbf{LSTM}
& \textbf{KI-SSDM}
& \textbf{KI Gain vs.\ Physics} \\
\midrule
\multirow{2}{=}{\textbf{Matched (S1)}}
& Planned
& 17.26 & 17.26 & 18.25 & 17.26 & 34.69
& \textbf{17.28}
& $-0.01$~h ($-0.1\%$) \\
& Probabilistic
& 19.33 & 19.32 & 19.97 & 19.33 & 35.22
& \textbf{19.46}
& $-0.14$~h ($-0.7\%$) \\
\midrule
\multirow{2}{=}{\textbf{Rate Discrepancy (S2)}}
& Planned
& 11.46 & 11.44 & 11.57 & 11.32 & 22.32
& \textbf{11.24}
& \textbf{+0.21~h (+1.9\%)} \\
& Probabilistic
& 13.46 & 13.40 & 13.80 & 13.35 & 22.34
& \textbf{13.15}
& \textbf{+0.30~h (+2.3\%)} \\
\midrule
\multirow{2}{=}{\textbf{Emission Distortion (S3)}}
& Planned
& 17.38 & 17.23 & 18.93 & 17.28 & 34.93
& \textbf{17.17}
& \textbf{+0.20~h (+1.2\%)} \\
& Probabilistic
& 19.25 & 19.18 & 20.72 & 19.25 & 34.32
& \textbf{19.24}
& \textbf{+0.01~h (+0.1\%)} \\
\midrule
\multirow{2}{=}{\textbf{Combined Mismatch (S4)}}
& Planned
& 11.15 & 11.12 & 11.77 & 10.96 & 21.69
& \textbf{10.91}
& \textbf{+0.24~h (+2.1\%)} \\
& Probabilistic
& 13.14 & 13.14 & 13.29 & 13.09 & 21.73
& \textbf{12.66}
& \textbf{+0.48~h (+3.6\%)} \\
\midrule
\multirow{2}{=}{\textbf{Weak Sensor Gap (S5)}}
& Planned
& 27.12 & 27.14 & 27.37 & 27.29 & 36.23
& \textbf{27.65}
& $-0.53$~h ($-2.0\%$) \\
& Probabilistic
& 27.43 & 27.43 & 27.59 & 27.60 & 38.68
& \textbf{28.00}
& $-0.56$~h ($-2.0\%$) \\
\midrule
\multirow{2}{=}{\textbf{Load Exponent Misspecification (S6)}}
& Planned
& 20.12 & 20.10 & 15.11 & 15.80 & 26.59
& \textbf{14.12}
& \textbf{+5.99~h (+29.8\%)} \\
& Probabilistic
& 20.61 & 20.60 & 17.50 & 18.19 & 26.65
& \textbf{17.21}
& \textbf{+3.40~h (+16.5\%)} \\
\bottomrule
\end{tabular}
\end{table*}
\paragraph{Robustness in the Null Regime (S1)}
When the physics model is fully specified (\texttt{fatigue\_matched}), an unconstrained learned correction risks overfitting. As shown in Table~\ref{tab:scenario_breakdown}, the LSTM baseline reaches $\mathrm{MAE} = 34.69$~h. KI-SSDM's correction is bounded, and in this scenario it changes forecasts little whichever candidate validation selects: KI-SSDM achieves \textbf{17.28~h} MAE, statistically indistinguishable from the nominal Physics SSM (17.26~h, a difference of just $-0.01$~h, paired $t$-test $p = 0.86$).

\paragraph{Structural Recovery under Load Misspecification (S6) and Gain Concentration}
A critical evaluation for knowledge-informed hybrid prognostics is whether the learned discrepancy can compensate for structural physical parameter misspecification without destabilising training. In Scenario S6, the true fatigue life exponent is $p = 4.5$, whereas the nominal physics model assumes $p = 3.0$. Under this structural mismatch, purely physics-driven Physics SSM and Bayesian PF incur large forecasting errors ($\mathrm{MAE} = 20.12$~h and $20.10$~h, respectively). KI-SSDM successfully compensates for the missing load-stress dependency via its candidate hypothesis pool (with 3 fleets selecting Basis RBF and 2 fleets selecting Deep tanh with stacking weights $w \in [0.75, 1.0]$), driving MAE down to \textbf{14.12~h}---achieving a \textbf{+5.99~h (+29.8\%) prognostic gain} over uncorrected physics (paired $t$-test $p = 0.0056$, Holm-adjusted $p = 0.0277$, winning 5 out of 5 fleets). Against the other learned corrections the S6 differences are not significant: Basis Hybrid (15.80~h, +1.68~h, $p_t = 0.40$) and Inverse PINN (15.11~h, +0.98~h, $p_t = 0.13$). Against Deep LSTM (26.59~h) the gain is +12.47~h ($p_t = 0.032$, $p_{\mathrm{Holm}} = 0.096$).

Crucially, an arithmetic audit of Table~\ref{tab:scenario_breakdown} reveals that the aggregate prognostic gain over pure physics is heavily concentrated in this misspecification regime. Across the six scenarios, the planned-duty total gain is $+6.11$~h (sum of rounded table values is $+6.10$~h), of which Scenario S6 contributes 98.2\% (+5.99~h). Under probabilistic duty, the total gain across scenarios is $+3.49$~h, of which S6 contributes 97.4\% (+3.40~h). Outside S6, the mean gain against Physics SSM across non-S6 regimes is +0.02~h Planned (+0.018~h Probabilistic), with slight advantages in S2 (+0.21~h), S3 (+0.20~h), and S4 (+0.24~h), and slight conservative differences in S1 ($-0.01$~h) and S5 ($-0.53$~h). The empirical value of the learned discrepancy therefore lies in correcting structural rate misspecification; in the other scenarios the bounded correction and validation selection keep KI-SSDM close to uncorrected physics. Small 0.01~h variations between displayed gains and direct differences reflect rounding of underlying floating-point values.

\paragraph{Resistance to Telemetry Blackouts (S5)}
Under extended data dropouts (\texttt{fatigue\_weak\_gap}), the LSTM baseline exhibits elevated prediction errors (36.23~h MAE Planned, 38.68~h Probabilistic). KI-SSDM reaches 27.65~h (planned) and 28.00~h (probabilistic), slightly worse than Physics SSM (27.12~h and 27.43~h; $-0.53$~h and $-0.56$~h); telemetry dropouts are the one regime in which the learned correction did not help. In this scenario the 120--240~h blackout is combined with sensor noise about 3.6 times higher than elsewhere (0.16 against 0.045), so the validation data carry too little information to fit the S2 rate discrepancy that is also present.

\paragraph{Telemetry Distortion and Combined Mismatches (S3--S4)}
In Scenario S3 (unmodeled nonlinear emission distortion), purely data-driven sequence architectures degrade to 34.93~h MAE, and Inverse PINN drifts to 18.93~h. KI-SSDM achieves \textbf{17.17~h} Planned MAE (+1.76~h over Inverse PINN, $+9.3\%$, $p_t = 0.041$, $p_{\mathrm{Holm}} = 0.16$), consistent with the GLS nuisance marginalisation limiting the effect of emission distortion on damage tracking. In Scenario S4 (simultaneous wear acceleration and emission distortion), KI-SSDM achieves \textbf{10.91~h} Planned MAE (+0.24~h gain vs.\ Physics SSM and +0.86~h gain vs.\ Inverse PINN).

\begin{figure*}[pos=!t]
  \centering
  \includegraphics[width=0.96\textwidth, keepaspectratio]{Fig2_paired_gain_forest.png}
  \caption{Paired prognostic gain forest plot ($\Delta\MAE = \MAE_{\mathrm{Comparator}} - \MAE_{\mathrm{KI\text{-}SSDM}}$) with 95\% paired confidence intervals across six degradation regimes and both duty protocols (30 fleets). Positive values indicate superior accuracy for KI-SSDM; LSTM gains (+13.01~h planned, +11.53~h probabilistic) are omitted from the plot and reported in Table~\ref{tab:overall_results}.}
  \label{fig:gain_forest}
\end{figure*}

\subsection{Conformal Calibration and Coverage Validation}
\label{subsec:conformal_results}

Table~\ref{tab:conformal_results} evaluates uncertainty calibration across test units across all 30 benchmark fleets. KI-SSDM's origin-stratified split-conformal prediction intervals attain an empirical test coverage of \textbf{93.3\%} [92.0\%, 94.6\%] under Planned Duty and \textbf{92.5\%} [90.9\%, 94.0\%] under Probabilistic Duty (fleet-level 95\% $t$ intervals, $n = 30$), above the nominal 90\% level ($1 - \alpha = 0.90$). Because every method is wrapped in the same split-conformal calibration, all six reach 90.3--93.4\% coverage, and the comparison rests on \emph{sharpness} at equal nominal coverage. KI-SSDM's mean interval widths are \textbf{82.48~h} (Planned) and \textbf{96.39~h} (Probabilistic), about 2.8 and 2.4 times narrower than Deep LSTM (229.20~h and 232.37~h) and similar to Basis Hybrid (81.47~h and 97.32~h) and Inverse PINN (81.83~h and 97.20~h). Its Winkler interval scores ($\mathrm{IS}_{0.10} = 94.20$ Planned, $110.04$ Probabilistic) are close to the best, Basis Hybrid's 94.01 (Planned), and the best under Probabilistic duty.


\begin{table}[!t]
\centering
\caption{Uncertainty quantification and conformal
coverage across 30 fleets (nominal coverage:
$1-\alpha=0.90$). Pointwise coverage denotes
marginal coverage at inspection origins;
simultaneous coverage requires coverage
at all origins within each unit.}
\label{tab:conformal_results}

\begingroup
\renewcommand{\arraystretch}{1.15}
\setlength{\tabcolsep}{1.5pt}
\scriptsize

\begin{tabularx}{\columnwidth}{
@{}
>{\raggedright\arraybackslash}X
>{\centering\arraybackslash}p{0.145\columnwidth}
>{\centering\arraybackslash}p{0.115\columnwidth}
>{\centering\arraybackslash}p{0.115\columnwidth}
>{\centering\arraybackslash}p{0.105\columnwidth}
>{\centering\arraybackslash}p{0.115\columnwidth}
@{}
}
\toprule
\textbf{Method}
& \textbf{Policy}
& \shortstack{\textbf{Point}\\\textbf{Cov.}}
& \shortstack{\textbf{Sim.}\\\textbf{Cov.}}
& \shortstack{\textbf{Width}\\\textbf{(h)}}
& \shortstack{\textbf{Interval}\\\textbf{Score}}\\
\midrule

\textbf{KI-SSDM (Proposed)}
& Planned
& \textbf{93.3\%}
& \textbf{79.0\%}
& \textbf{82.48}
& \textbf{94.20}\\

& Probabilistic
& \textbf{92.5\%}
& \textbf{79.2\%}
& \textbf{96.39}
& \textbf{110.04}\\
\midrule

\textbf{Basis Hybrid}
& Planned
& 92.9\%
& 77.7\%
& 81.47
& 94.01\\

& Probabilistic
& 93.0\%
& 80.0\%
& 97.32
& 110.79\\
\midrule

\textbf{Inverse PINN}
& Planned
& 90.3\%
& 74.4\%
& 81.83
& 97.87\\

& Probabilistic
& 90.9\%
& 75.8\%
& 97.20
& 113.75\\
\midrule

\textbf{Bayesian PF}
& Planned
& 93.0\%
& 79.6\%
& 88.47
& 100.05\\

& Probabilistic
& 92.8\%
& 79.8\%
& 102.78
& 115.85\\
\midrule

\textbf{Physics SSM}
& Planned
& 93.1\%
& 79.2\%
& 88.58
& 100.13\\

& Probabilistic
& 92.7\%
& 79.6\%
& 102.39
& 115.33\\
\midrule

\textbf{Deep LSTM}
& Planned
& 93.4\%
& 82.9\%
& 229.20
& 256.46\\

& Probabilistic
& 92.2\%
& 82.3\%
& 232.37
& 265.86\\

\bottomrule
\end{tabularx}
\endgroup
\end{table}


\begin{figure*}[pos=!t]
  \centering
  \includegraphics[width=0.95\textwidth, keepaspectratio]{Fig3_calibration_sharpness.png}
  \caption{Calibration and sharpness trade-off across 30 fleets (480 test units). Top row: Empirical test coverage versus nominal 90\% coverage target, (KI-SSDM: $93.3\%$ Planned, $92.5\%$ Probabilistic). Bottom row: mean prediction-interval width (sharpness); KI-SSDM bands average $82.5$~h against $229.2$~h for Deep LSTM (Planned).}
  \label{fig:cal_sharpness}
\end{figure*}

\subsection{Methodological Ablation Study: Residual Dynamics Impact}
\label{subsec:ablation_study}

To isolate the direct empirical contribution of the learned residual discrepancy dynamics within KI-SSDM, an architectural ablation study was evaluated across all 30 benchmark fleets (Table~\ref{tab:ablation_results}). The ablation compares the full proposed framework (Residual Dynamics ON: active bounded discrepancy $\rho(\mathbf{u}, D)$, Kalman nuisance tracking, and validation linear pooling) against the ablated baseline (Residual Dynamics OFF). Residual OFF evaluates the same fitted KI-SSDM, under the same posterior and random draws, with the discrepancy switched off ($b = 0$); it is a forecast-only ablation and not a refitted physics model, which is why it differs from Physics SSM (Table~\ref{tab:overall_results}). Because the backbone parameters were estimated jointly with the discrepancy, removing $\rho$ leaves a backbone that was not fitted to stand alone, and Residual OFF is less accurate than the separately fitted Physics SSM.

The results demonstrate:
\begin{enumerate}[label=(\alph*),leftmargin=2.2em]
    \item \textbf{Planned Duty Performance}: With residual dynamics enabled, the Planned fleet-mean MAE was \textbf{16.40~h} against $18.21$~h without, a mean gain of \textbf{+1.81~h} that did not reach significance (95\% CI $[-0.004, +3.63]$~h, paired $t$-test $p = 0.050$, Wilcoxon $p = 0.092$; 12 of 19 decisive fleets).
    \item \textbf{Probabilistic Duty Performance}: Under Probabilistic duty, residual dynamics achieves a statistically significant improvement, reducing MAE from $20.02$~h to \textbf{18.29~h} (\textbf{+1.74~h gain}, 95\% CI $[+0.37, +3.10]$~h, paired $t$-test $p = 0.015$, Wilcoxon $p = 0.0019$, 16 of 22 decisive fleets).
\end{enumerate}


\begin{table}[!t]
\centering
\caption{Ablation of learned residual dynamics
across 30 benchmark fleets (480 test units).}
\label{tab:ablation_results}

\begingroup
\renewcommand{\arraystretch}{1.18}
\setlength{\tabcolsep}{1.5pt}
\scriptsize

\begin{tabularx}{\columnwidth}{
@{}
>{\raggedright\arraybackslash}X
>{\centering\arraybackslash}p{0.115\columnwidth}
>{\centering\arraybackslash}p{0.115\columnwidth}
>{\centering\arraybackslash}p{0.225\columnwidth}
>{\centering\arraybackslash}p{0.095\columnwidth}
>{\centering\arraybackslash}p{0.13\columnwidth}
@{}
}
\toprule
\textbf{Duty Policy}
& \shortstack{\textbf{Residual}\\\textbf{ON (h)}}
& \shortstack{\textbf{Residual}\\\textbf{OFF (h)}}
& \shortstack{\textbf{Gain (h)}\\\textbf{[95\% CI]}}
& \shortstack{\textbf{Fleet}\\\textbf{W/L}}
& \shortstack{$\boldsymbol{p_t}$ /\\$\boldsymbol{p_W}$}\\
\midrule

Planned Duty
& \textbf{16.40}
& 18.21
& \shortstack{\textbf{+1.81}\\$[-0.004,\,+3.63]$}
& 12/7
& \shortstack{0.050 /\\0.092}\\[2pt]

Probabilistic Duty
& \textbf{18.29}
& 20.02
& \shortstack{\textbf{+1.74}\\$[+0.37,\,+3.10]$}
& 16/6
& \shortstack{0.015 /\\0.0019}\\

\bottomrule
\end{tabularx}
\endgroup
\end{table}


\subsection{Test-Bench Operating Characterisation}
\label{subsec:bench_characterisation}

The physical test bench (Fig.~\ref{fig:bsfd_setup}) was evaluated across four traversal feed speeds (5,000--20,000~mm/min, corresponding to servomotor speeds of 500--2,000~rpm) under unladen carriage ($m_{\mathrm{tbl}} = 25.0$~kg) and 150~kg dynamic payload ($M_{\mathrm{tot}} = 175.0$~kg) conditions. Table~\ref{tab:bench} reports steady-speed torque RMS and mean nut temperature recorded directly from controller-resident OPC~UA telemetry. Because the physical rig was not run to destructive fatigue failure, these experimental measurements characterise operating loads and thermal sensitivity rather than remaining life; no uncomputed run-to-failure RUL is claimed for the bench.


\begin{table}[!t]
\centering
\caption{Measured steady-speed torque RMS,
mean temperature, and ISO~3408-5 relative
rated life under unladen and 150~kg loading.
Life is normalized to the unladen
10,000~mm/min condition (1.00).}
\label{tab:bench}

\begingroup
\renewcommand{\arraystretch}{1.15}
\setlength{\tabcolsep}{2pt}
\footnotesize

\begin{tabularx}{\columnwidth}{
@{}
>{\raggedright\arraybackslash}X
*{4}{>{\centering\arraybackslash}p{0.125\columnwidth}}
@{}
}
\toprule
\textbf{Parameter}
& \textbf{5,000}
& \textbf{10,000}
& \textbf{15,000}
& \textbf{20,000}\\
\midrule

Torque RMS, unladen
(N$\cdot$m)
& 1.30 & 1.76 & 2.39 & 3.06\\

Torque RMS, 150~kg
(N$\cdot$m)
& 1.26 & 1.88 & 2.49 & 3.24\\

Mean temperature,
unladen ($^\circ$C)
& 31.1 & 36.6 & 49.3 & 56.1\\

Mean temperature,
150~kg ($^\circ$C)
& 46.2 & 47.3 & 55.3 & 67.1\\

Relative rated life,
unladen
& 1.65 & 1.00 & 0.76 & 0.67\\

Relative rated life,
150~kg
& 1.67 & 1.00 & 0.76 & 0.65\\

\bottomrule
\end{tabularx}

\par\vspace{2pt}
{\scriptsize\raggedright
\textit{Note:} Column headings indicate
feed speed in mm/min.\par}

\endgroup
\end{table}


Two physical characteristics are revealed by the telemetry. First, the preloaded double-nut assembly is dominated by its 1.2~kN internal preload: the 150~kg external payload, being supported by the recirculating linear guideway rails, increases the per-nut equivalent dynamic axial load by at most 1.2\% (1,216.6~N vs 1,202.1~N at 20,000~mm/min), so payload alters nominal fatigue life by only a few per cent while rotational speed modulates life markedly (Table~\ref{tab:bench}). Second, the operating torque level is heavily influenced by thermal expansion: in same-day experimental traversals, torque RMS decreased by 30.7\%, 38.7\%, and 45.4\% from cold start to warm steady state at 5,000, 10,000, and 15,000~mm/min as the ball nut warmed by 8--17~$^\circ$C. Consequently, torque-based condition indicators must decouple thermal fluctuations from structural wear, directly motivating the continuous-time Ornstein--Uhlenbeck nuisance process $B(t)$ formulated in Section~\ref{subsec:nuisance_b}. Absolute rated life additionally requires the catalogue dynamic load rating~\cite{PMI_Catalog,ISO_3408_5}.

\begin{figure*}[pos=!t]
  \centering
  \includegraphics[width=0.96\textwidth, keepaspectratio]{fig_realtime_speed_load.png}
  \caption{Real-time edge-deployed operating-regime RUL sensitivity and horizon projections across traversal feed speeds (5,000--20,000~mm/min) for unladen ($m_{\mathrm{tbl}} = 25.0$~kg) and 150~kg carriage payload ($M_{\mathrm{tot}} = 175.0$~kg) regimes on the industrial BSFD platform. (a) Unladen carriage configuration demonstrating speed-dependent kinematic life compression from $8{,}271$~h ($5{,}000$~mm/min = 500~rpm) to $8{,}132$~h ($20{,}000$~mm/min = 2,000~rpm). (b) 150~kg payload configuration illustrating accelerated degradation due to elevated dynamic axial loading and thermal dissipation ($8{,}143$~h down to $8{,}056$~h). Shaded error bars depict 90\% conformal prediction intervals for KI-SSDM and 90\% uncertainty intervals for comparators against the nominal ISO~3408-5 design life reference horizon ($8{,}260$~h nominal). Note: In figure panels, PI-DSSDM denotes the proposed KI-SSDM framework.}
  \label{fig:realtime_speed_load}
\end{figure*}

To evaluate edge deployment feasibility on real industrial hardware, the online filtering and RUL projection pipeline was deployed to process streaming OPC~UA telemetry directly. As detailed in Section~\ref{subsec:realtime_features}, raw OPC~UA telemetry arrives at an asynchronous cadence of $\approx 100$~ms ($\approx 10$~samples/s, median $9.2$~Hz) and is upsampled to $100$~samples/s ($10$~ms grid) over $1.0$~s sliding analysis windows ($500$~ms hop). On the edge computing host, the algorithmic computation time for the deterministic state update and feature extraction executes in under $1$~ms per window, consuming less than $0.2\%$ of the $500$~ms window-hop budget and demonstrating real-time edge viability without controller communication bottlenecks. Across operating speeds, Figure~\ref{fig:realtime_speed_load} visualises the projected operational life horizons:
\begin{enumerate}[label=(\alph*),leftmargin=2.2em]
    \item \textbf{Guideway Load Decoupling}: Because the carriage and workpiece mass is borne by the recirculating linear guideway rails, adding 150~kg of dynamic payload introduces only a modest increment in equivalent screw-nut axial stress (at most 1.2\% under ISO~3408-5). Consequently, the projected RUL shifts by only a few tens of hours between unladen and laden operation (e.g., $8{,}271 \to 8{,}143$~h at 5,000~mm/min, and $8{,}132 \to 8{,}056$~h at 20,000~mm/min), corroborating mechanical load decoupling.
    \item \textbf{Speed-Dependent Kinematic Dissipation}: Across all evaluated methods, higher traversal velocities increase cyclic reversal frequencies and thermo-mechanical friction dissipation ($67.1^\circ\mathrm{C}$ at 20,000~mm/min), leading to monotonic life compression from $8{,}271$~h at 5,000~mm/min down to $8{,}056$~h at 20,000~mm/min.
    \item \textbf{Monotonicity and Conformal Sharpness}: While unconstrained sequence models (Deep LSTM) exhibit localized non-monotonic sensitivity across feedrate steps due to unconstrained recurrent gating, proposed KI-SSDM strictly enforces thermodynamic wear monotonicity ($dD/dt \ge 0$) and maintains sharp calibrated conformal prediction intervals ($211.2$~h width) across all operational regimes.
\end{enumerate}
\subsection{Downstream Decision Support: Risk-Constrained Maintenance Scheduling and Lifecycle Cost}
\label{subsec:decision_support}

In modern engineering informatics and intelligent manufacturing asset management, condition-based prognostics must provide actionable decision support rather than isolated scalar predictions~\cite{Aivaliotis_2021, D_He_2025}. A pervasive deficiency in current physics-informed deep learning literature is that prognostics terminates at evaluating empirical RUL point metrics or theoretical error distributions, without translating predictive uncertainty into concrete operational maintenance policies. In real-world factory environments, machine tool asset managers face asymmetric operational risks: catastrophic in-service tool collisions during critical finishing cuts incur severe financial penalties, whereas overly conservative, premature component overhauls discard valuable remaining equipment life.

\subsubsection{Lifecycle Maintenance Cost Model under Asymmetric Risk}
To evaluate the downstream decision utility of KI-SSDM, we formulate an industrial lifecycle maintenance cost model. Let $t_0$ denote the current inspection time, and let $T \in \R^+$ denote the ground-truth remaining useful life until critical fatigue failure ($D(t) = \Dcrit = 1.0$). An asset manager must schedule an overhaul lead time $\tau \ge 0$, committing to take the feed drive offline for overhaul at planned maintenance epoch $t_{\mathrm{sched}} = t_0 + \tau$.

The operational outcome and associated economic cost $\mathcal{C}(\tau; T)$ per maintenance cycle depend on whether the planned overhaul precedes or succeeds the true physical failure:

\begin{equation}
\label{eq:cost_model}
\mathcal{C}(\tau;T)=
\begin{cases}
C_p+c_w\max(0,T-\tau),
& \tau\leq T,\\[3pt]
C_p+C_u,
& \tau>T.
\end{cases}
\end{equation}

The first case represents a planned overhaul ($\tau\leq T$), whereas the second corresponds to catastrophic failure ($\tau>T$). Here, $C_p$ denotes the baseline cost of a scheduled, planned replacement (normalized to $C_p = 1.0$), including replacement ball-screw procurement, planned technician shifts, and scheduled cell downtime; $C_u$ represents the severe catastrophic failure penalty (including ruined high-value workpieces, spindle tool-holder impact damage, emergency overtime labour, and unplanned production line stoppage; we take $r_c = C_u / C_p = 8$ as a representative assumption and test 4 and 16 in Table~\ref{tab:cost_sens}); and $c_w = C_p / L_{\mathrm{design}} = 1.0 \times 10^{-4}$~h$^{-1}$ (a declared assumption, one replacement cost per $L_{\mathrm{design}} = 10{,}000$ nominal hours) represents the marginal amortized cost of discarded, unconsumed useful component life.


\subsubsection{Conformal Risk-Constrained Decision Policy}
Standard heuristic maintenance strategies typically schedule interventions based on the point estimate: $\tau_{\mathrm{point}} = \widehat{\mathrm{RUL}}$. Under symmetric error distributions, this strategy inherently implies a 50\% chance of scheduled overhaul occurring after true failure ($\mathbb{P}(\tau_{\mathrm{point}} > T) \approx 0.50$). Even with slightly conservative estimation biases, empirical point-estimate scheduling yields high catastrophic failure rates. Heuristic Bayesian credible intervals ($\mathrm{RUL}_{0.05}^{\mathrm{PF}}$) attempt to mitigate this risk, but particle impoverishment and observation misspecification compromise their empirical coverage guarantees.

In contrast, KI-SSDM directly connects its origin-stratified split-conformal calibration to the decision process. By selecting the scheduled maintenance lead time as the lower conformal bound:
\begin{equation}
    \tau^*(\alpha) = \mathrm{RUL}_{\mathrm{lo}}(t_0) = \max\bigl\{0,\ \widehat{\mathrm{RUL}}(t_0) - \hat{q}_{1-\alpha}(t_0)\bigr\},
    \label{eq:conformal_lead_time}
\end{equation}
the finite-sample marginal coverage guarantee (Eq.~\eqref{eq:coverage_guarantee}), which holds under exchangeability of calibration and test units, bounds the probability of in-service breakdown:

\begin{equation}
\label{eq:conformal_risk_bound}
\begin{aligned}
\mathbb{P}(\text{Unplanned Failure})
&= \mathbb{P}\bigl(\tau^*(\alpha)>T\bigr)\\
&= \mathbb{P}\bigl(
T<\mathrm{RUL}_{\mathrm{lo}}(t_0)
\bigr)\\
&\leq \alpha.
\end{aligned}
\end{equation}

For a risk budget of $\alpha = 0.10$, the marginal probability of in-service failure is therefore at most 10\%, and the narrower the interval, the less useful life is given up to achieve it.

\subsubsection{Comparative Lifecycle Cost Evaluation Across 480 Fleet Units}
Table~\ref{tab:decision_maintenance_costs} evaluates the economic performance of KI-SSDM's risk-constrained maintenance policy against representative industrial maintenance strategies evaluated across all 480 held-out test units across the 30 benchmark fleets (1,776 unit-origin maintenance decisions per duty mode evaluated on surviving units). Periodic maintenance is scheduled at the 10th-percentile calibration-fleet failure age ($L_{10}$ age per fleet).

\begin{table}[!t]
\centering
\caption{Comparative evaluation of maintenance decision policies across 480 held-out test units under Planned Duty ($C_p = 1.0$, $C_u = 8.0$, $c_w = 1.0 \times 10^{-4}$~h$^{-1}$, 1,776 decisions). Periodic replacement is set to the fleet-specific calibration $L_{10}$ age.}
\label{tab:decision_maintenance_costs}
\begingroup
\renewcommand{\arraystretch}{1.15}
\setlength{\tabcolsep}{2pt}
\scriptsize
\begin{tabularx}{\columnwidth}{
@{}
>{\raggedright\arraybackslash}X
>{\centering\arraybackslash}p{0.13\columnwidth}
>{\centering\arraybackslash}p{0.14\columnwidth}
>{\centering\arraybackslash}p{0.14\columnwidth}
>{\centering\arraybackslash}p{0.17\columnwidth}
@{}
}
\toprule
\textbf{Maintenance Strategy}
& \shortstack{\textbf{Failure}\\\textbf{Rate}}
& \shortstack{\textbf{Discarded}\\\textbf{Life}}
& \shortstack{\textbf{Expected}\\\textbf{Cost} ($\bm{C_p}$)}
& \shortstack{\textbf{Saving vs.}\\\textbf{Periodic} ($\bm{L_{10}}$)}\\
\midrule
Corrective (Run-to-Failure)
& 100.0\% & 0.0\% & 9.00 & $-347.0\%$ \\
Periodic (Calibration $L_{10}$ Age)
& 12.6\% & 24.2\% & 2.01 & $0.0\%$ (Benchmark) \\
\midrule
\multicolumn{5}{@{}l}{\textit{Point-Estimate Decision Policies ($\tau = \widehat{\mathrm{RUL}}$):}} \\
KI-SSDM Point Estimate
& 45.6\% & 3.7\% & 4.65 & $-130.7\%$ \\
Inverse PINN Point Estimate
& 46.4\% & 4.0\% & 4.71 & $-134.1\%$ \\
Deep LSTM Point Estimate
& 51.5\% & 7.0\% & 5.12 & $-154.2\%$ \\
\midrule
\multicolumn{5}{@{}l}{\textit{Conformal Risk-Constrained Decision Policies ($\tau = \mathrm{RUL}_{\mathrm{lo}}$):}} \\
Inverse PINN Conformal Lower
& 4.8\% & 9.1\% & 1.39 & $+31.1\%$ \\
Bayesian PF Conformal Lower
& 4.3\% & 9.4\% & 1.35 & $+33.2\%$ \\
Physics SSM Conformal Lower
& 4.2\% & 9.4\% & 1.34 & $+33.4\%$ \\
Basis Hybrid Conformal Lower
& 4.1\% & 8.9\% & 1.33 & $+33.8\%$ \\
Deep LSTM Conformal Lower
& 3.9\% & 18.1\% & 1.32 & $+34.3\%$ \\
\textbf{KI-SSDM Conformal Lower ($\bm{\alpha=0.10}$)}
& \textbf{3.9\%} & \textbf{9.0\%} & \textbf{1.31} & \textbf{+34.7\%} \\
\bottomrule
\end{tabularx}
\endgroup
\end{table}

As shown in Table~\ref{tab:decision_maintenance_costs}:
\begin{enumerate}[label=(\alph*),leftmargin=2.2em]
    \item \textbf{Failure Rate Suppression}: Periodic replacement at the calibration $L_{10}$ benchmark incurs an empirical catastrophic failure rate of $12.6\%$ (above the nominal 10\%; each fleet's $L_{10}$ age is estimated from its 24 calibration units), while discarding an average of $24.2\%$ of usable component life on robust units ($\mathbb{E}[\mathcal{C}] = 2.01\,C_p$). Uncalibrated point-estimate policies (KI-SSDM point: $45.6\%$; Inverse PINN: $46.4\%$; LSTM: $51.5\%$) fail to account for prognostic uncertainty, suffering high catastrophic failure rates due to symmetric error distributions around the point forecast and incurring severe economic penalties ($\mathbb{E}[\mathcal{C}] = 4.65\text{--}5.12\,C_p$, representing a $-130\%$ to $-154\%$ penalty relative to periodic replacement).
    \item \textbf{Conformal Risk Bounding}: Scheduling overhauls at the split-conformal lower bound ($\tau^* = \mathrm{RUL}_{\mathrm{lo}}$) restricts failure rates to at most $4.8\%$ under planned duty and $5.1\%$ under probabilistic duty across all evaluated methods ($3.9\%$ and $4.7\%$ for KI-SSDM). This is within the risk budget $\alpha = 0.10$ guaranteed by Eq.~\eqref{eq:conformal_risk_bound}. With the symmetric absolute-error score, the lower tail alone is bounded by $\alpha/2 = 5\%$ only if prediction errors are symmetric; empirically all lower-bound policies stayed at or below 5.1\%.
    \item \textbf{Lifecycle Cost Reduction}: Every prognostic method saves 30--35\% when scheduling at its conformal lower bound rather than relying on point forecasts or periodic overhaul. Under Planned Duty, KI-SSDM has the lowest expected cost ($\mathbb{E}[\mathcal{C}] = 1.31\,C_p$, a \textbf{34.7\% reduction relative to periodic $L_{10}$ replacement}) while discarding $9.0\%$ of usable life; Deep LSTM is close ($1.32\,C_p$, $34.3\%$) but discards twice as much life ($18.1\%$) because of its wider intervals (Table~\ref{tab:conformal_results}). Under Probabilistic Duty, KI-SSDM reaches $1.38\,C_p$ (\textbf{31.6\% saving}, $9.6\%$ discarded life); here Deep LSTM is cheaper ($1.32\,C_p$) because its failure rate is lower ($3.9\%$ vs $4.7\%$), at the price of discarding $17.0\%$ of usable life. With $c_w = C_p/L_{\mathrm{design}}$ applied to the nominal benchmark clock, discarded life costs little (about $0.003\,C_p$ per decision), so expected cost is driven almost entirely by the failure rate.
\end{enumerate}

Table~\ref{tab:cost_sens} repeats the evaluation for $C_u/C_p \in \{4, 8, 16\}$ and for a fleet-scale price of discarded life, $c_w = C_p/\bar{T}_{\mathrm{CAL}}$, where $\bar{T}_{\mathrm{CAL}}$ is the mean calibration-unit life of each fleet. Scheduling at the KI-SSDM conformal lower bound remains the cheapest planned-duty policy, or tied cheapest, in all six settings, with savings of 23--47\% over periodic $L_{10}$ replacement; with discarded life priced at the fleet scale, Deep LSTM becomes the most expensive planned-duty policy ($1.51$ vs $1.40\,C_p$ at $C_u/C_p = 8$). Under probabilistic duty KI-SSDM, Basis Hybrid and Bayesian PF are within $0.02\,C_p$ of one another, and Deep LSTM, whose failure rate is lowest (3.9\%), is cheapest in four of the six settings.

\begin{table*}[pos=!t]
\centering
\caption{Sensitivity of expected cost per decision ($C_p = 1$) to the failure-cost ratio $C_u/C_p$ and to the price of discarded life $c_w$ (manuscript value $10^{-4}$~h$^{-1}$, or fleet-scale $C_p/\bar{T}_{\mathrm{CAL}}$). Conformal-lower-bound policies; 1,776 decisions per duty. Bold: lowest cost in the row.}
\label{tab:cost_sens}
\renewcommand{\arraystretch}{1.12}
\setlength{\tabcolsep}{4pt}
\small
\begin{tabular}{@{}lccccccccc@{}}
\toprule
\textbf{Duty} & $\bm{C_u/C_p}$ & $\bm{c_w}$ & \textbf{Periodic $L_{10}$} & \textbf{KI-SSDM} & \textbf{Basis Hybrid} & \textbf{Physics SSM} & \textbf{Bayesian PF} & \textbf{Inverse PINN} & \textbf{Deep LSTM} \\
\midrule
Planned & 4 & $10^{-4}$ & 1.51 & \textbf{1.16} & 1.17 & 1.17 & 1.17 & 1.19 & \textbf{1.16} \\
Planned & 4 & fleet & 1.74 & \textbf{1.25} & 1.26 & 1.27 & 1.27 & 1.28 & 1.35 \\
Planned & 8 & $10^{-4}$ & 2.01 & \textbf{1.31} & 1.33 & 1.34 & 1.35 & 1.39 & 1.32 \\
Planned & 8 & fleet & 2.25 & \textbf{1.40} & 1.42 & 1.44 & 1.44 & 1.48 & 1.51 \\
Planned & 16 & $10^{-4}$ & 3.02 & \textbf{1.63} & 1.66 & 1.68 & 1.69 & 1.77 & 1.64 \\
Planned & 16 & fleet & 3.25 & \textbf{1.71} & 1.75 & 1.77 & 1.78 & 1.86 & 1.82 \\
\midrule
Probabilistic & 4 & $10^{-4}$ & 1.51 & 1.19 & 1.19 & 1.20 & 1.19 & 1.21 & \textbf{1.16} \\
Probabilistic & 4 & fleet & 1.74 & 1.29 & \textbf{1.28} & 1.30 & 1.29 & 1.30 & 1.34 \\
Probabilistic & 8 & $10^{-4}$ & 2.01 & 1.38 & 1.37 & 1.39 & 1.38 & 1.41 & \textbf{1.32} \\
Probabilistic & 8 & fleet & 2.25 & \textbf{1.47} & \textbf{1.47} & 1.49 & 1.48 & 1.51 & 1.49 \\
Probabilistic & 16 & $10^{-4}$ & 3.02 & 1.75 & 1.74 & 1.78 & 1.76 & 1.82 & \textbf{1.64} \\
Probabilistic & 16 & fleet & 3.25 & 1.85 & 1.84 & 1.88 & 1.86 & 1.92 & \textbf{1.81} \\
\bottomrule
\end{tabular}
\end{table*}

% ========================  7. DISCUSSION  ========================
\section{Discussion}
\label{sec:discussion}

\subsection{Physical Interpretability and Identifiability}
\label{subsec:disc_identifiability}

A central vulnerability of hybrid neural systems is the mathematical confounding between physical model parameters and learned neural weights. In unconstrained models, neural networks can easily compensate for erroneous physical constants, destroying interpretability.

In KI-SSDM, identifiability is structurally protected through three orthogonal mechanisms:
\begin{enumerate}[label=(\roman*),leftmargin=2.2em]
    \item \textbf{Separation of Timescales and Scale Constraints}: The OU parameter $\theta_B$ (simulated value $1/18\,\mathrm{h}^{-1}$) governs reversible operational fluctuations with characteristic half-life $t_{1/2} = \ln(2)/\theta_B \approx 12.5$~h, whereas irreversible fatigue $D(t)$ evolves monotonically over hundreds of hours in the benchmark. In addition, explicit scale constraints (unit stationary variance of $B$, so that its amplitude is carried by the bounded loadings $\gamma_j \in [-2, 2]$, and zero-mean frailty $\mathbb{E}[\eta_u] = 0$) prevent confounding between nuisance amplitude and wear scaling.
    \item \textbf{Authority Clamping}: The strict bound on $\rho(t) \in [-b, b]$ restricts neural authority to $[e^{-b}, e^b]$, ensuring strictly positive wear multipliers ($\alpha > 0$) and preventing the network from absorbing primary scale parameters.
    \item \textbf{GLS Conditioning}: Emission parameters $\boldsymbol{\beta}$ are resolved conditionally via convex quadratic optimisation, ensuring unique global solutions given the design matrix.
\end{enumerate}

\subsection{Computational Cost and Deployment Outlook}
\label{subsec:edge_deployment}

Offline training of the full 30-fleet benchmark (six methods, 2,220 units) took 18.8~h of summed stage time (about 22~h end-to-end wall clock) on 12 CPU worker processes in MATLAB. Online, KI-SSDM needs per window one damage-state update, a conditional Gaussian update of the nuisance state for each frailty support point, and a forward pass of the bounded discrepancy network; this deterministic step executes in $<1$~ms per 500~ms window hop, consuming less than $0.2\%$ of available cycle time. RUL forecasts are obtained by forward simulation of damage paths and need only be refreshed at inspection origins. Formal multi-core hardware-in-the-loop benchmarking on dedicated industrial micro-controllers remains future work.
\subsection{Limitations and Scope}
\label{subsec:limitations}

While KI-SSDM represents a rigorous step forward, its boundaries must be noted:
\begin{enumerate}[label=(\alph*),leftmargin=2.2em]
    \item \textbf{Single Failure Family Scope}: The formulation focuses on rolling-contact fatigue. Competing failure modes, such as catastrophic seal tearing or winding breakdown, require multi-hazard competing-risk extensions.
    \item \textbf{Fleet Calibration}: While parameterised from physical PMI catalogue ratings and rig measurements, full confirmatory claims will benefit from multi-year industrial fleet tracking across diverse machining centres.
    \item \textbf{Single-Axis Kinematic Focus and Multi-Axis Couplings}: The present validation focuses on a single horizontal feed axis. In 5-axis machine tools, dynamic tilting moments across guideways and non-uniform thermal expansion gradients along the screw shaft ($x$-dependent pitch drift) introduce additional geometric couplings. Extending the state-space formulation to multi-axis kinematic chains with coupled thermo-mechanical digital twins represents a fruitful research avenue.
    \item \textbf{Empirical Evidence Gate and Multiplicity Bounds}: In the synthetic benchmark hierarchy, while KI-SSDM demonstrates decisive, statistically significant gains under severe structural load-exponent misspecification (Scenario S6: $+5.99$~h gain, $p = 0.0056$, Holm-adjusted $p = 0.0277$), under rate discrepancy (Scenario S2: $+0.21$~h gain) and telemetry dropouts (Scenario S5: $-0.53$~h), the framework performs at parity with uncorrected physics. The reproducibility package applies a stricter pre-declared evidence check that requires the multiplicity-adjusted lower confidence bound of every scenario-by-comparator gain to be positive; only 40 of these 60 contrasts pass (failing lower bounds range from $-0.41$ to $-9.74$~h), and the checks on numerical resolution of late-origin contrasts and on final-objective convergence of every saved checkpoint are not all met. Furthermore, the benchmark is dimensionless, with a simulated mean lifetime of $371.9$~h on a nominal time scale, and is not calibrated to the physical screw's dynamic load rating; its three telemetry channels are abstract and sampled every 8~h. Finally, the physical test bench was used to characterise operating loads and thermal sensitivity from controller telemetry alone; run-to-failure validation of KI-SSDM, and measurement of its latency on embedded hardware, require further experiments.
\end{enumerate}

% ========================  8. CONCLUSIONS  ========================
\section{Conclusions}
\label{sec:conclusions}

This paper introduced KI-SSDM, a knowledge-informed state-space degradation modelling framework for predictive maintenance and risk-constrained decision support in precision mechatronics. By explicitly formalizing engineering domain knowledge into structured state-space elements---uniting continuous ISO~3408-5 fatigue mechanics~\cite{ISO_3408_5}, dynamic tribology, thermodynamic dissipation, and unit frailty---the framework isolates irreversible wear from operational nuisance. Analytical Generalised Least Squares (GLS) marginalisation and Kalman innovations evaluate exact Gaussian telemetry likelihoods in closed-form $\mathcal{O}(K)$ time without particle degeneracy. Deep neural discrepancies parameterised by smooth $C^\infty$ tanh-bounded skip connections capture structural misspecifications under strict physical authority bounds, ensuring strictly positive wear multipliers ($\alpha > 0$) and monotonic degradation ($\mathrm{d}D/\mathrm{d}t \ge 0$). Furthermore, validation-weighted linear pooling prevents empirical divergence, origin-stratified split-conformal prediction provides finite-sample coverage guarantees, and an integrated decision layer optimizes lifecycle maintenance costs.

Across 30 independent synthetic fleets (2,220 simulated units; 480 held-out test units; six scenarios), KI-SSDM reached fleet-mean MAEs of 16.40~h (planned duty) and 18.29~h (probabilistic duty). After Holm adjustment it was significantly more accurate than Inverse PINN (planned duty, $p_{\mathrm{Holm}} = 0.0079$) and Deep LSTM (both duties, $p_{\mathrm{Holm}} < 10^{-9}$), and not significantly different from Physics SSM, Bayesian PF ($p_{\mathrm{Holm}} = 0.107$) and Basis Hybrid ($p_{\mathrm{Holm}} = 0.418$). Its advantage over uncorrected physics was concentrated under load-exponent misspecification (+5.99~h, +29.8\%, planned duty), with a mean gain of +0.02~h in the matched regimes. Calibrated intervals covered 93.3\% and 92.5\% of test cases at mean widths of 82.5~h and 96.4~h, while unit-simultaneous coverage was 79\%. Scheduling at the conformal lower bound kept the failure rate at 3.9\% (planned) and 4.7\% (probabilistic) and reduced expected cost by 34.7\% (planned) and 31.6\% (probabilistic) relative to periodic $L_{10}$ replacement; comparators scheduled the same way saved 30--35\%, so the calibrated lower bound, rather than the point forecast, drives the saving. Future work will validate the framework on bench run-to-failure data, measure embedded deployment latency, and extend it to multi-axis coupling and competing failure mechanisms.


% ========================  DECLARATIONS & ACKNOWLEDGEMENTS  ========================

\section*{CRediT Authorship Contribution Statement}
\textbf{Harikrishna S. Thota}: Conceptualization, Methodology, Software, Formal analysis, Validation, Visualization, Writing -- original draft, Writing -- review \& editing. \\
\textbf{P.S. Saikrishna}: Investigation, Validation, Data curation, Supervision, Writing -- review \& editing. \\
\textbf{Nagahanumaiah}: Supervision, Project administration, Resources, Funding acquisition.

\section*{Declaration of Competing Interest}
The authors declare that they have no known competing financial interests or personal relationships that could have appeared to influence the work reported in this paper.

\section*{Data and Code Availability}
The synthetic benchmark generator, model and baseline implementations, evaluation scripts, and all result tables are archived in the supplementary reproducibility package (\texttt{PI\_DSSDM\_Reproducibility\_Package\_v2}). The package is provided to reviewers as supplementary material, and the complete code repository will be made publicly available upon publication at \url{https://github.com/harithota/KI-SSDM-BSFD-Prognostics}. The code runs in MATLAB (tested on R2022b and R2025b); the Deep Learning Toolbox is required for the LSTM baseline, and the Parallel Computing Toolbox is used if available. Running the package script \texttt{main\_reproduce\_all} regenerates every table and figure from the archived run outputs; a fresh 30-fleet run takes 18.8~h of summed stage time (about 22~h end-to-end wall clock) on 12 CPU workers. Test-bench telemetry is available from the corresponding author on reasonable request.

\section*{Funding}
This work was funded by the Ministry of Heavy Industries (MHI), Government of India, under the SMDDC [DPP1701203] project.

\section*{Acknowledgements}
The authors gratefully acknowledge the institutional computing support and experimental test-bench facilities provided by the Central Manufacturing Technology Institute (CMTI), Bengaluru, as well as collaborative insights from the Indian Institute of Technology Tirupati.

\section*{Declaration of Generative AI in the Writing Process}
During the preparation of this work, the authors utilised Google Antigravity / Gemini for language editing, \LaTeX{} syntax formatting, and stylistic refinement of mathematical notation. After utilising this tool, the authors reviewed and edited the content as needed and take full responsibility for the content of the published article.

% ========================  REFERENCES  ========================
\bibliographystyle{unsrtnat}
\bibliography{References_uploaded_selected_v2_corrected}

\end{document}

